\chapter{Multiple Fourier Series}
\label{ch:c7}

The purpose of this chapter is to give a brief introduction to several parts of the theory of multiple Fourier series. While from the abstract point of view this represents only a particular instance of Fourier analysis on compact abelian groups, our emphasis is on the connection between the analysis on the $n$-torus\index{n-torus@n-torus} and Euclidean $n$-dimensional space. Thus, in Section 2 we consider the process of ``periodization\index{periods@periods}'' applied to functions on $E_n$, which gives us the Poisson summation formula, and in Section 3 we consider the corresponding periodization for multiplier operators. In the fourth and fifth sections we consider another topic, the summability of multiple Fourier series. We study that subject at this stage, instead of earlier in the nonperiodic context (in which very similar results could be proved) mainly for reasons of elegance and appropriateness. The results that are presented are valid only in the case $n > 1$, and are obtained by techniques different from those originally developed for treating similar problems in the more familiar one-dimensional case.

\section{Elementary Properties}

Let $\Lambda$ denote the additive group of points in $E_n$ having integral coordinates (the addition being, of course, the one derived from the vector structure of $E_n$). $\Lambda$ will be called the unit lattice\index{unit lattice@unit lattice}. We consider the coset space $E_{n^1}/\Lambda$ and we identify, in the usual manner, functions defined on $E_n/\Lambda$ with periodic functions on $E_n$. More precisely, a function $f$ satisfying $f(x + m) = f(x)$ whenever $m \in \Lambda$ is identified with the function whose value at the coset determined by $x$ is $f(x)$. The elements of $\Lambda$ are the periods of these functions.

There exists a natural identification of $E_n/\Lambda$ with the $n$-torus $T_n = \{(e^{2\pi i x_1}, e^{2\pi i x_2}, \ldots, e^{2\pi i x_n}) \in C^n: (x_1, x_2, \ldots, x_n) \in E_n\}$. This identification is given by the mapping $(x_1, x_2, \ldots, x_n) \to (e^{2\pi i x_1}, e^{2\pi i x_2}, \ldots, e^{2\pi i x_n})$ and, from it, we obtain the standard identification of periodic functions on $E_n$ with functions on the $n$-torus.

A set $D \subset E_n$ will be called a fundamental domain\index{fundamental domain@fundamental domain} if every point of $E_n$ has exactly one translate with respect to $\Lambda$ in $D$. A periodic function on $E_n$ is, clearly, uniquely determined by its restriction to a fundamental domain. Obviously, there are infinitely many fundamental domains; one that is

both simple and appropriate for our purposes is the fundamental cube\index{fundamental cube@fundamental cube} $Q_n = \{x \in E_n: -1/2 \leq x_j < 1/2, j = 1, 2, \ldots, n\}$. Integration on $T_n$ can be defined in terms of the Lebesgue integral on $Q_n$: if $f$ is a function on $T_n$ we write

\begin{equation}\label{eq:c7-1-1}
\int_{T_{n}}f dx=\int_{Q_{n}}f dx
\end{equation}

and interpret this equality in the following way. Any function $f$ on $T_n$, as we have seen, gives rise to a periodic one on $E_n$, and the restriction of this last function to $Q_n$ is the one appearing in the right side of \eqref{eq:c7-1-1}. The function $f$ on $T_n$ is defined to be measurable if the corresponding function on $Q_n$ is Lebesgue measurable; if the latter is integrable then the integral on the left side of \eqref{eq:c7-1-1} is defined to be equal to the integral on the right side. In the same way, the $L^p$ spaces on $T_n$ are identified with the $L^p$ spaces on $Q_n$. Similarly, we can identify the class of finite Borel measures on $T_n$ with the class of finite Borel measures concentrated on the fundamental cube $Q_n$. We shall use the notation $L^p(T_n)$ and $\mathcal{B}(T_n)$ for the $L^p$ spaces on $T_n$ and the Borel measures on $T_n$, respectively. It should be pointed out, however, that the class of continuous functions on $T_n$, $C(T_n)$, does not correspond to the class of continuous functions on $Q_n$, but only to those functions which remain continuous on $E_n$ when extended by periodicity. $C(T_n)$ is a Banach subspace of $L^\infty(T_n)$ (that is, we endow the former space with the $L^\infty$ norm). We record here the elementary inclusion relations

\[
\mathcal{B}(T_{n})\supset L^{p}(T_{n})\supset C(T_{n})
\]

and remark that, as usual, $L^{1}(T_{n})$ may be considered as identified with the subspace of absolutely continuous measures of $\mathcal{B}(T_{n})$.

We shall now discuss Fourier series of measures and functions in $T_n$. To every element $\mu \in \mathcal{B}(T_n)$ we assign the Fourier series

\begin{equation}\label{eq:c7-1-2}
d\mu\sim\sum_{m\in\Lambda}a_{m}e^{2\pi i m\cdot x},
\end{equation}

where

\begin{equation}\label{eq:c7-1-3}
a_{m}=\int\limits_{T_{n}}e^{-2\pi i m\cdot x}d\mu(x)
\end{equation}

are the Fourier-Stieltjes coefficients of $d\mu$.

To facilitate formal manipulation with Fourier coefficients it is convenient to consider, for each measure $\mu$ in $\mathcal{B}(T_{n})$, the continuous linear functional

\begin{equation}\label{eq:c7-1-4}
{\cal L}_{\mu}\colon f\to\intop_{T_{n}}f d\mu,
\end{equation}

which it defines on the space $C(T_n)$. By the Riesz-Markov representation theorem we know that, conversely, every continuous linear functional on $C(T_n)$ is of the form \eqref{eq:c7-1-4} for some $\mu \in \mathcal{B}(T_n)$ with $\|L_\mu\| = \|d_\mu\|$, where the norm on the left is that of the linear functional $L_\mu$ and the norm on the right is the total measure of $\mu$. This identification makes it easy to define the convolution of two measures\index{convolution@convolution!of two measures (on n-torus)@of two measures (on $n$-torus)} $\mu_1$ and $\mu_2$ in $\mathcal{B}(T_n)$: it is the measure $\mu$ which represents (via \eqref{eq:c7-1-4}) the linear functional\index{multiplication formula@multiplication formula!for two measures on n-torus@for two measures on n-torus}

\begin{equation}\label{eq:c7-1-5}
\dot{f}\to\intop_{T_{n}}f d\mu=\intop_{T_{n}}\intop_{T_{n}}f(x+y)d\mu_{1}(x)d\mu_{2}(y).
\end{equation}

From this definition it is clear that $\|d\mu\| \leq \|d\mu_1\| \|d\mu_2\|$ and, with this operation of convolution, $\mathcal{B}(T_n)$ is a commutative Banach algebra. If we put $f(x) = e^{-2\pi i m \cdot x}$ in \eqref{eq:c7-1-5} we obtain the multiplication formula

\begin{equation}\label{eq:c7-1-6}
d\mu_{1}*d\mu_{2}\sim\sum_{m\in\Lambda}a_{m}b_{m}e^{2\pi i m\cdot x},
\end{equation}

where $\{a_m\}$ and $\{b_m\}$ are the Fourier-Stieltjes coefficients of $d\mu_1$ and $d\mu_2$ respectively.

A cursory examination of \eqref{eq:c7-1-5} also leads to an alternative definition of $d\mu_1*d\mu_2$. It is the measure $\mu$ with the property that

\[
\mu(E)=\int\limits_{T_{n}}\mu_{1}(E-y)d\mu_{2}(y)
\]

for any Borel set $E$. From this it is clear that if $\mu_1$ (or $\mu_2$) is absolutely continuous then $\mu$ is also absolutely continuous. The measure $\mu_1$, however, is absolutely continuous if and only if it has a Radon--Nikodym derivative $f\in L^1(T_n)$. In this case the Radon--Nikodym derivative of $\mu$ is the function $h$ defined by

\[
h(x)=\int\limits_{T_{n}}f(x-y)d\mu_{2}(y).
\]

This integral converges absolutely for almost every $x$ by virtue of Fubini's theorem. If $\mu_2$ is also absolutely continuous, with Radon--Nikodym derivative $g \in L^1(T_n)$, then

\[
h(x)=\int\limits_{T_{n}}f(x-y)g(y)d y.
\]

Thus, we see that $L^{1}(T_{n})$ inherits the convolution structure from $\mathcal{B}(T_{n})$. In particular, the multiplication formula \eqref{eq:c7-1-6} holds for $L^{1}$ functions as well.

In order to formulate the basic facts concerning completeness of the Fourier expansion, it is convenient to refer to any finite sum of the form $\sum a_m e^{2\pi im \cdot x}$ as a trigonometric polynomial\index{trigonometric polynomial@trigonometric polynomial}.

\begin{theorem}
\label{thm:c7-1-7}
(i) The trigonometric polynomials are dense in $C(T_n)$ and $L^p(T_n), 1 \leq p < \infty$.
(ii) Suppose that for some $\mu \in \mathcal{B}(T_n)$, $\int_{T_n} e^{-2\pi i m \cdot x} \, d\mu(x) = 0$ for all $m \in \Lambda$; then $\mu = 0$.
(iii) Suppose $f \in L^2(T_n)$ and $f \sim \sum_{m \in \Lambda} a_m e^{2\pi i m \cdot x}$, then
\[
\sum_{m\in\Delta}|a_{m}|^{2}=\|f\|_{2}^{2}.
\]

\end{theorem}

The correspondence $f \leftrightarrow \{a_m\} = \{\int_{T_n} e^{-2\tau_{nm} \cdot \gamma} f(x) \, dx\}$ is a unitary mapping of $L^2(T_n)$ onto $l^2(\Lambda)$.

\textsc{Proof.}\ Since the trigonometric polynomials form an algebra which separates points of the $n$-torus, which contains the constants and is closed under complex conjugation, we can apply the Stone-Weierstrass theorem in order to obtain the density of the trigonometric polynomials in $C(T_n)$. This approximation property together with the density of $C(T_n)$ in $L^p(T_n)$, $1 \leq p < \infty$, implies the corresponding approximation property in $L^p(T_n)$.

In order to show (ii) we observe that the property $\int_{T_n} e^{-2\pi i m \cdot x} \, d\mu(x) = 0$ for all $m \in \Lambda$ is equivalent to

\[
\int\limits_{T_{n}}P(x)d\mu(x)=0
\]

for each trigonometric polynomial $P$. But, by (i), this implies that $\int_{T_n} f(x) \, d\mu(x) = 0$ for all $f \in C(T_n)$. It follows, therefore, that $\mu = 0$.

Given $f\in L^{2}(T_{n})$ then, if $N$ is a positive integer, the infimum of the expression $\|f-\sum_{|m|\leq N}a_{m}e^{2\pi im\cdot x}\|_{2}$ is attained when the numbers $a_{n}$ are the Fourier coefficients\index{Fourier coefficients@Fourier coefficients!of a function on n-torus@of a function on $n$-torus} $\int_{T_{n}}e^{-2\pi im\cdot x}f(x)\,dx$ of $f$. This assertion follows, as is well-known, from the fact that the functions $\{e^{2\pi im\cdot x}\}$ are mutually orthogonal and normalized (in the $L^{2}$ norm). Because each $f\in L^{2}(T_{n})$ can be approximated by trigonometric polynomials we know, therefore, that $\lim_{N\to\infty}\|f-\sum_{|m|\leq N}a_{m}e^{2\pi im\cdot x}\|_{2}=0$. Consequently,

\[
\left\|f\right\|_{2}-\left\|\sum_{|m|\leq N}a_{m} e^{2\pi i m\cdot x}\right\|_{2}\rightarrow0
\]

as $N \to \infty$ and so $\|f\|_2 = (\sum_{m \in \Lambda} |a_m|^2)^{1/2}$. This shows that the mapping $f \to \{a_m\}$ is an isometry. If it is also onto, then this mapping must be unitary. In fact, let $\{a_m\}$ be given with $\sum_{m \in \Lambda} |a_m|^2 < \infty$ and put $s_N(x) = \sum_{|m| \leq N} a_m e^{2\pi i m \cdot x}$. Clearly, $\|s_{N_1} - s_{N_2}\|_2 = (\sum_{N_1 < |m| \leq N_2} |a_m|^2)^{1/2}$, if $N_1 < N_2$. Therefore, $\{s_N\}$ is a Cauchy sequence in $L^2(T_n)$ and, thus, it converges.

in the norm to some function $f \in L^{2}(T_{n})$. Since

\[
a_{m}=\intop_{\dot{T}_{n}}s_{N}(x)e^{-2\pi i x\cdot m}\,d x\qquad{\text{if }}|m|\,\leq\,N,
\]

it follows that

\[
a_{m}=\int\limits_{T_{n}}f(x)e^{-2\pi i x\cdot m}d x\qquad\text{for all }m\in\Lambda.
\]

This shows (iii) and, therefore, the theorem is proved.

We now derive some useful corollaries of this result.

\begin{corollary}
\label{cor:c7-1-8}
Suppose $f \in L^1(T_n)$ and $\sum_{m \in \Lambda} |a_m| < \infty$, where $\{a_m\}$ are the Fourier coefficients of $f$; then $f$ can be modified on a set of measure zero so that it is in $C(T_n)$ and equals $\sum_{m \in \Lambda} a_m e^{2\pi im \cdot x}$ for all $x \in T_n$.
\end{corollary}

\textsc{Proof.}\ It follows from the hypotheses that the function whose values are $f(x) - \sum_{m \in \Lambda} a_m e^{2\pi im \cdot x}$ has zero Fourier coefficients; thus, by (ii) of Theorem \ref{thm:c7-1-7}, this function must be 0 almost everywhere.

When $k$ is a positive integer the class $C^{(k)}(Q_n) (= C^{(k)}(T_n))$ consists of the restrictions to $Q_n$ of all periodic functions on $E_n$ belonging to the class $C^{(k)} = C^{(k)}(E_n).\footnote{That is, those periodic functions $f$ having everywhere defined continuous derivatives $D^\alpha f$ whenever $\alpha = (\alpha_1, \alpha_2, \ldots, \alpha_n)$ is an $n$-tuple of nonnegative integers satisfying $\alpha_1 + \alpha_2 + \cdots + \alpha_n \leq k$. Recall that in Chapter I the operator $D^\alpha$ was defined to be 
\[
\frac{\partial^{\alpha_{1}}}{\partial x_{1}^{\alpha_{1}}}\frac{\partial^{\alpha_{2}}}{\partial x_{2}^{\alpha_{2}}}\cdot\cdot\cdot\frac{\partial^{\alpha_{n}}}{\partial x_{n}^{\alpha_{n}}}.
\]}$

\begin{corollary}
\label{cor:c7-1-9}
Suppose $f \in C^{(k)}(T_n)$ for $k > n/2$, then $\sum_{m \in \Lambda} |a_m| < \infty$, where $\{a_m\}$ are the Fourier coefficients of $f$.
\end{corollary}

\textsc{Proof.}\ In the first chapter we defined $x^{\alpha}$ for $x = (x_1, x_2, \ldots, x_n) \in E_n$ and $\alpha = (\alpha_1, \alpha_2, \ldots, \alpha_n)$ an $n$-tuple of nonnegative integers to be the number $x_{1}^{a_1} x_{2}^{a_2} \cdots x_{n}^{a_n}$ (we also adopted the convention that $0^0 = 1$). Using this notation and integrating by parts we obtain

\[
\int_{T_{n}}(D^{\alpha}f)(x)e^{-2\pi i m\cdot x}\mathop{}{d}x=(2\pi i m)^{\alpha}\int_{T_{n}}f(x)e^{-2\pi i m\cdot x}\mathop{}{d}x=(2\pi i m)^{\alpha}a_{m},
\]

whenever $f \in C^{(k)}(T_n)$ and $\alpha_1 + \alpha_2 + \cdots + \alpha_n \leq k$. Since $D^\alpha f$ is continuous

We warn the reader that, in that chapter, the sum $\alpha_1 + \alpha_2 + \cdots + \alpha_n$ was denoted by $|\alpha|$; however, in the proof of Theorem \ref{thm:c7-1-7} we used the symbol $|m|$ to denote the Euclidean norm of the lattice point $m$. In order to avoid confusion we will not use here the notation $|\alpha|$ introduced in the first chapter.

it must belong to $L^{2}(T_{n})$; therefore, by Theorem \ref{thm:c7-1-7} (iii)

\begin{equation}\label{eq:c7-1-10}
\sum_{\alpha_{1}+\cdot+\cdot+\alpha_{n}=k}\left\{\sum_{m\in\Lambda}|a_{m}|^{2}[(2\pi m)^{\alpha}]^{2}\right\}<\infty.
\end{equation}

Moreover, a simple calculation shows that there exists a constant $c = c(k, n)$, depending only on the dimension n and on k, such that

\[
\sum_{\alpha_{1}+\cdot\cdot\cdot+\alpha_{n}=k}[(2\pi m)^{\alpha}]^{2}\geq c|m|^{2k}.
\]

Thus, by Schwarz's inequality

\[
\begin{aligned}\sum_{|m|>0}|a_{m}|\leq&\sum_{|m|>0}|a_{m}|\biggl(\sum_{\alpha_{1}+\cdots+\alpha_{n}=k}[(2\pi m)^{\alpha}]^{2}\biggr)^{1/2}c^{-1/2}|m|^{-k}\\\leq&\biggl(\sum_{|m|>0}|a_{m}|^{2}\sum_{\alpha_{1}+\cdots+\alpha_{n}=k}[(2\pi m)^{\alpha}]^{2}\biggr)^{1/2}\biggl(\sum_{|m|>0}|m|^{-2k}\biggr)^{1/2}c^{-1/2}.\end{aligned}
\]

When $k > n/2$ the sum $\sum_{|m| > 0} |m|^{-2k}$ is finite; hence, the last expression must also be finite in this case by virtue of inequality \eqref{eq:c7-1-10}. Thus, the corollary is proved.

The following corollary is the familiar Riemann--Lebesgue theorem in the present context:

\begin{corollary}
\label{cor:c7-1-11}
If $f \in L^1(T_n)$ and $f \sim \sum_{m \in \Lambda} a_m e^{2\pi i m \cdot x}$ then $a_m \to 0$ as $|m| \to \infty$.
\end{corollary}

\textsc{Proof.}\ Suppose first that $f \in L^{2}(T_{n})$. Then the conclusion $a_{m} \to 0$ is an immediate consequence of the fact that $\sum_{m \in \Lambda} |a_{m}|^{2} < \infty$ (see Theorem \ref{thm:c7-1-7} part (iii)). For the general $f \in L^{1}(T_{n})$, given $\varepsilon > 0$, we can find $f_{1}$ and $f_{2}$ such that $f = f_{1} + f_{2}$ with $f_{1} \in L^{2}(T_{n})$ and $\|f_{2}\|_{1} < \varepsilon$. Then, if $\{a_{m}^{(j)}\}$ denote the Fourier coefficients of $f_{i}$, we have $\lim_{|m| \to \infty} a_{m}^{(1)} = 0$ and $|a_{m}^{(2)}| \leq \|f_{2}\|_{1} < \varepsilon$. Thus, $\limsup_{|m| \to \infty} |a_{m}| \leq \varepsilon$, where $a_{m} = a_{m}^{(1)} + a_{m}^{(2)}, m \in \Lambda$, are the Fourier coefficients of $f$. Since $\varepsilon > 0$ is arbitrary, this establishes the corollary.

\section{The Poisson Summation Formula}

Instead of continuing a step-by-step development of the analogy between the properties of Fourier series and those of Fourier integrals, we now turn to the question which is the main concern of this chapter. We can phrase it in the following general way: Suppose we are given an ``element'' of a function space associated with $E_n$ (e.g., a function on $E_n$, an operator on functions of $E_n$, etc.), what is its periodic analog? That is, what object corresponds to it on the $n$-torus $T_n$? In addition we shall be interested in understanding how we can deduce the properties of such an

object in its periodic form from those we have already established for its non-periodic form.

We consider these matters first in the context of functions defined on $E_{n}$. In order to obtain a better idea of the problems involved we proceed formally, disregarding for the moment all questions of convergence.

Let, then, $f$ be some (appropriate) function on $E_n$. There are at least two ways of obtaining a periodic function from $f$. The first construction is elementary and does not involve Fourier analysis. We simply write

\begin{equation}\label{eq:c7-2-1}
\sum_{m\in\Lambda}f(x+m).
\end{equation}

Since this (formal) sum is taken over all the lattice points of $\Lambda$ it is clearly periodic (for the passage from $x$ to $x + m'$ merely permutes the terms in \eqref{eq:c7-2-1}). We shall refer to the passage from $f$ to the sum \eqref{eq:c7-2-1} as the periodization\index{periodization@periodization!of functions@of functions} of $f$.

In order to describe the second approach we write

\begin{equation}\label{eq:c7-2-2}
f(x)\sim\int\limits_{E_{n}}\hat{f}(y)e^{2\pi i x\cdot y}d y,
\end{equation}

which is simply the Fourier inversion formula, where

\[
\hat{f}(y)=\int\limits_{E_{n}}f(x)e^{-2\pi i x\cdot y}d x.
\]

Then a periodic analog of the function $f$(given by \eqref{eq:c7-2-2}) is

\begin{equation}\label{eq:c7-2-3}
\sum_{m\in\Lambda}\hat{f}(m)e^{2\pi i x\cdot m}.
\end{equation}

The main thrust of the Poisson summation formula is to the effect that the two approaches to a periodic analog of f given by \eqref{eq:c7-2-1} and \eqref{eq:c7-2-3} are essentially identical. This conclusion can be formulated precisely in many ways. The simplest and most direct is the following.

\begin{theorem}
\label{thm:c7-2-4}
Suppose $f \in L^1(E_n)$. Then the series $\sum_{m \in \Lambda} f(x + m)$ converges in the norm of $L^1(Q_n) (= L^1(T_n))$. The resulting function in $L^1(T_n)$ has the Fourier expansion
\begin{equation}\label{eq:c7-2-5}
\sum_{m\in\Delta}\hat{f}(m)e^{2\pi i x\cdot m}.
\end{equation}

\end{theorem}

This means that $\{ \hat{f}(m) \}$ are the Fourier coefficients of the $L^{1}$ function defined by the series $\sum_{m \in \Lambda} f(x + m)$, where, for any $y \in E_n$,

\[
\hat{f}(y)=\int\limits_{E_{n}}f(x)e^{-2\pi i y\cdot x}d x.
\]

\textsc{Proof.}\ If $Q_{n}-m$ is the translate of $Q_{n}$ by the lattice point $m$ we then have

\[
\int\limits_{Q_{n}}\left|\sum_{m\in\Lambda}f(x+m)\right|d x\leq\sum_{m\in\Lambda}\int\limits_{Q_{n}}\left|f(x+m)\right|d x=\sum_{m\in\Lambda}\intop_{Q_{n}^{-m}}\left|f(x)\right|d x.
\]

Since the fundamental cube $Q_{n}$ is a fundamental domain, the translates $Q_{n} - m$ are mutually disjoint and their union is $E_{n}$. Thus,

\[
\sum_{m\in\Delta}\int\limits_{Q_{n}-m}\left|f(x)\right|dx=\int\limits_{E_{n}}\left|f(x)\right|dx<\infty.
\]

This shows that the series $\sum_{m \in \Lambda} f(x + m)$ converges (absolutely) in the norm of $L^1(Q_n)$. Using the same interchange of summation and integration we can evaluate the Fourier coefficients of $\sum_{m \in \Lambda} f(x + m)$. In fact,

\[
\begin{aligned}\int\limits_{Q_{n}}\Big(\sum_{m^{\prime}\in\Lambda}f(x+m^{\prime})\Big)e^{-2\pi i m\cdot x}dx&=\sum_{m^{\prime}\in\Lambda}\int\limits_{Q_{n}-m^{\prime}}f(x)e^{-2\pi i m\cdot x}dx\\&=\int\limits_{F_{n}}f(x)e^{-2\pi i n\cdot x}dx\\&=\hat{f}(m).\end{aligned}
\]

The argument given above for the $L'$ convergence also shows that this interchange is justified and. thus, the fact that $\sum_{m'\in\Lambda} f(x + m')$ has the Fourier expansion \eqref{eq:c7-2-5} is demonstrated.

The following corollary is a useful special case of Theorem \ref{thm:c7-2-4}:

\begin{corollary}
\label{cor:c7-2-6}
Suppose $\hat{f}(y) = \int_{E_n} f(x) e^{-2\pi i x \cdot y} \, dx$ and
\[
f(x)=\int\limits_{E_{n}}\hat{f}(y)e^{2\pi i x\cdot y}dy
\]

with $|f(x)| \leq A(1 + |x|)^{-n-\delta}$ and $|\hat{f}(y)| \leq A(1 + |y|)^{-n-\delta}$, $\delta > 0$ (thus, both \textit{and} $\hat{f}$ can be assumed to be continuous). Then,
\begin{equation}\label{eq:c7-2-7}
\sum_{m\in\Lambda}f(x+m)=\sum_{m\in\Lambda}\hat{f}(m)e^{2\pi i m\cdot x},
\end{equation}

and, in particular,
\begin{equation}\label{eq:c7-2-8}
\sum_{m\in\Delta}f(m)=\sum_{m\in\Delta}\hat{f}(m).
\end{equation}

\end{corollary}

The four series in \eqref{eq:c7-2-7} and \eqref{eq:c7-2-8} converge absolutely. \footnote{The identity \eqref{eq:c7-2-8} is referred to as the Poisson summation formula, but we shall also use the same name for \eqref{eq:c7-2-7} and, more generally, for Theorem \ref{thm:c7-2-4}.}

\textsc{Proof.}\ Because of our assumption on $f$, the Fourier series

\[
\sum_{m\in\Lambda}\hat{f}(m)e^{2\pi i m\cdot x}
\]

converges absolutely. Therefore, according to Corollary \ref{cor:c7-1-8} and Theorem \ref{thm:c7-2-4}, the function $\sum_{m\in\Lambda} f(x + m)$ can be modified on a set of measure zero so as to equal the continuous function $\sum_{m\in\Lambda} \hat{f}(m)e^{2\pi im\cdot x}$ everywhere. By comparison with $\sum_{m\in\Lambda} (1 + |m|)^{-n-\delta}$, however, we see that $\sum_{m\in\Lambda} f(x + m)$ is a uniformly convergent series whose terms are continuous functions. Thus, its sum is already continuous and we obtain \eqref{eq:c7-2-7} for every $x$.

We shall now give two illustrations of how the Poisson summation\index{Poisson summation formula@Poisson summation formula} formula may be used. In the first example we consider the analog of the Fourier inversion problem considered in Chapter I; that is, we ask to what extent can the series

\[
\sum_{m\in\Lambda}a_{m}e^{2\pi i m\cdot x}
\]

be ``summed'' to $f(x)$ when $a_m = \int_{T_n} f(x) e^{-2\pi im \cdot x} dx$ with $f \in L^1(T_n)$. In view of what we learned in the first chapter it is natural to try to replace the above sum by the limit

\begin{equation}\label{eq:c7-2-9}
\lim_{\varepsilon\to0}\sum_{m\in\Lambda}\Phi(\varepsilon m)a_{m}e^{2\pi i m\cdot x},
\end{equation}

where $\Phi$ is an appropriate continuous function with $\Phi(0) = 1$. Motivated by Corollary \eqref{cor:c7-2-6} we make the following assumptions on $\Phi$:

\begin{equation}\label{eq:c7-2-10}
\left\{\begin{array}{l l}{\mathrm{(i)}}&{\Phi(y)=\hat{\varphi}(y)\qquad\text{with}\int_{E_{n}}\varphi(x)d x=1}\\ {\text{(ii)}}&{|\Phi(y)|\leq A(1+|y|)^{-n-\delta},|\varphi(x)|\leq A(1+|x|)^{-n-\delta},}\\ &{\quad\text{for some }\delta}\end{array}\right.
\end{equation}

\[
\delta>0.
\]

\begin{theorem}
\label{thm:c7-2-11}
Suppose $\Phi$ satisfies the conditions \eqref{eq:c7-2-10} above, $f \in L^{p}(T_{n})$ and $f(x) \sim \sum_{m \in \Lambda} a_{m} e^{2\pi i m \cdot x}$. Then
(a) If $1 \leq p < \infty$, the limit \eqref{eq:c7-2-9} converges to f in the norm of $L^p(T_n)$,
(b) If $f \in C(T_n)$, the limit \eqref{eq:c7-2-9} converges to $f$ uniformly;
(c) For any $f \in L^{1}(T_{n})$ the limit \eqref{eq:c7-2-9} converges to $f(x)$ at each point $x$ of the Lebesgue set of $f$ (and, thus, almost everywhere).
\end{theorem}

\textsc{Proof.}\ Let us write $\varphi_{\varepsilon}(x) = \varepsilon^{-n}\varphi(x/\varepsilon)$ and $\Phi^{\varepsilon}(y) = \Phi(\varepsilon y)$, $\varepsilon > 0$. Then we know that $\Phi^{\varepsilon} = (\varphi_{\varepsilon})^{\wedge}$. Moreover, because of \eqref{eq:c7-2-10} $\varphi_{\varepsilon}$ and $\Phi^{\varepsilon}$ satisfy the conditions of corollary \eqref{cor:c7-2-6} and, therefore,

\begin{equation}\label{eq:c7-2-12}
\sum_{m\in\Lambda}\Phi(\varepsilon m)e^{2\pi i m\cdot x}=\sum_{m\in\Lambda}\varphi_{\varepsilon}(x+m).
\end{equation}

Let us denote the right side of \eqref{eq:c7-2-12} by $K_{\varepsilon}(x)$. Then

\[
(f* K_{\varepsilon})(x)=\sum_{m\in\Lambda}\Phi(\varepsilon m)a_{m}e^{2\pi i m\cdot x},
\]

where the convolution is the one introduced in the first section (see, in particular, \eqref{eq:c7-1-6}). Because of our assumptions on $\Phi$, this last series converges absolutely. On the other hand,

\[
\intop_{T_{n}}\left|K_{\varepsilon}(x)\right|d x\leq\intop_{Q_{n}}\sum_{m\in\Lambda}\left|\varphi_{\varepsilon}(x+m)\right|d x=\intop_{E_{n}}\left|\varphi_{\varepsilon}(x)\right|d x=\intop_{E_{n}}\left|\varphi(x)\right|d x.
\]

Since $\|f*K_\varepsilon\|_p \leq \|f\|_p\|K_\varepsilon\|_1$ (where the norms are those of the spaces $L^p(T_n), 1 \leq p \leq \infty$), we see that the mappings

\[
M_{\varepsilon}\colon f\to\sum_{m\in\Lambda}\Phi(\varepsilon m)a_{m}e^{2\pi i x\cdot m}
\]

are uniformly bounded in $\varepsilon$, $\varepsilon > 0$, as operators on $L^p(T_n)$. By the continuity of $\Phi$ and the fact that $\int_{E_n} \varphi(x) \, dx = 1$ we have $\lim_{\varepsilon \to 0} \Phi(\varepsilon m) = 1$; therefore, $M_\varepsilon f \to f$ as $\varepsilon \to 0$ (in the $L^p$ norm) whenever $f$ is a trigonometric polynomial. The fact that the trigonometric polynomials are dense in $L^p(T_n)$, for $p < \infty$, and in $C(T_n)$, together with the uniform boundedness we just established, proves parts (a) and (b) of the theorem.

Before proving (c) we should first explain that we say $x$ is a point of\index{Lebesgue set@Lebesgue set!point of@point of} the \textit{Lebesgue set} of $f$ provided, after $f$ is continued periodically, $x$ is a point of the \textit{Lebesgue set} of the continued periodic function on $E_n$. It is clear that after an appropriate translation we may take our point $x$ to be the origin (which is the center of the fundamental cube $Q_n$). Let $\tilde{f}(x) = f(x)$ for $x \in Q_n$ and $\tilde{f}(x) = 0$ for $x \in E_n - Q_n$. Since the property of being a member of the \textit{Lebesgue set} of a function is a local one, $0$ is a point of the \textit{Lebesgue set} of $\tilde{f}$. Now,

\[
\begin{aligned}\sum_{m\in\Lambda}\Phi(\varepsilon m)a_{m}&=\int_{Q_{n}}f(x)K_\varepsilon(-x)dx=\sum_{m\in\Lambda}\int_{Q_{n}}f(x)\varphi_\varepsilon(-x+m)dx\\&=\int_{Q_{n}}f(x)\varphi_\varepsilon(-x)dx+\sum_{m\neq0}\int_{Q_{n}}f(x)\varphi_\varepsilon(-x+m)dx.\end{aligned}
\]

Because of our assumption on $\varphi$,

\[
\begin{aligned}|\varphi_{\varepsilon}(-x+m)|&\leq A\bigg(1+\bigg|\frac{-x+m}{\varepsilon}\bigg|\bigg)^{-n-\delta}\varepsilon^{-n}\\&=\varepsilon^{\delta}A(\varepsilon+|-x+m|)^{-(n+\delta)}\\&\leq\varepsilon^{\delta}A^{\prime}|m|^{-(n+\delta)},\quad\text{if}x\in Q_{n}\text{and}|m|\geq1.\end{aligned}
\]

Therefore,

\[
\left|\sum_{m\neq0}\int\limits_{Q_{n}}f(x)\varphi_{\varepsilon}(-x+m)d x\right|\leq\varepsilon^{\delta}A^{\prime \prime}\int\limits_{Q_{n}}\left|f(x)\right|d x\rightarrow0\qquad\text{ as }\varepsilon\rightarrow0.
\]

Also, $\int_{Q_n} f(x)\varphi_\varepsilon(-x)\,dx = \int_{E_n} \tilde{f}(x)\varphi_\varepsilon(-x)\,dx$. But the condition $|\varphi(x)| \leq A(1 + |x|)^{-n-\delta}$ allows us to apply Theorem \ref{thm:c1-1-25} of Chapter I to this last integral and we obtain

\[
\lim_{\varepsilon\to0}\int\limits_{E_{n}}\tilde{f}(x)\varphi_{\varepsilon}(-x)d x=\tilde{f}(0)=f(0).
\]

This proves part (c) and, thus, Theorem \ref{thm:c7-2-11} is completely proved.

Two particular special cases of this theorem are worth noting separately. First consider $\varphi(x) = c_n(1 + |x|^2)^{-(n+1)/2}$; then $\Phi(y) = e^{-2\pi|y|}$ (see Theorem \ref{thm:c1-1-14} of Chapter I). In this case the series occurring in \eqref{eq:c7-2-9} becomes (if we write $t$ instead of $\varepsilon$)

\begin{equation}\label{eq:c7-2-13}
\sum_{m\in\Lambda}e^{-2\pi\left|m\right|t}a_{m}e^{2\pi i m\cdot x}.
\end{equation}

This series, which is absolutely convergent for t > 0, is called the Poisson (or Abel--Poisson) integral of f; it is the convolution of f with the Poisson kernel

\[
P_{t}(x)=\sum_{m\in\Lambda}e^{-2\pi|m|t}e^{2\pi i m\cdot x}.
\]

It is a consequence of \eqref{eq:c7-2-12} that $P_t(x) \geq 0$ for $t > 0$ and $x \in Q_n$. Moreover,

\[
\intop_{Q_{n}}P_{t}(x)\mathop{}{d}x=\intop_{Q_{n}}1\mathop{}{d}x+\intop_{Q_{n}}\sum_{m\neq0}e^{-2\pi|m|t}e^{2\pi i m\cdot x}\mathop{}{d}x=1+0=1.
\]

The second example arises when $\Phi(y) = (1 - |y|^2)^{\alpha}$ if $|y| \leq 1$ and $\Phi(y) = 0$ if $|y| \geq 1$. By Theorem \ref{thm:c4-4-15} of Chapter IV we then have

\[
\begin{array}{r}{\varphi(x)=\Phi(x)=\pi^{-\alpha}\Gamma(\alpha+1)|x|^{-(n/2)-\alpha}J_{(n/2)+\alpha}(2\pi|x|).}\end{array}
\]

Thus, if $\alpha > (n - 1)/2$, it follows from Lemma \ref{lem:c4-3-11} of Chapter IV that $\Phi$ (and $\varphi$) satisfy hypotheses \eqref{eq:c7-2-10}. Therefore, in this case we obtain the appropriate convergence for the\index{Poisson integral@Poisson integral!for the n-torus@for the $n$-torus} $Riesz$ means\footnote{The number $(n - 1)/2$, the lower bound for the values $\alpha$ for which the functions $\Phi$ and $\varphi = \hat{\Phi}$ satisfy \eqref{eq:c7-2-10}, is called the critical index (for the convergence of the Riesz means\index{Riesz means@Riesz means} \eqref{eq:c7-2-14}).}\index{Poisson kernel@Poisson kernel!for the n-torus@for the $n$-torus}

\begin{equation}\label{eq:c7-2-14}
\lim_{R\to\infty}\sum_{|m|<R}\left(1\;-\;\frac{|m|^{2}}{R^{2}}\right)^{\alpha}a_{m}e^{2\pi i m\cdot x},\qquad\alpha>\frac{n-1}{2}.
\end{equation}

We summarize these facts as follows

\begin{corollary}
\label{cor:c7-2-15}
The conclusions of Theorem \ref{thm:c7-2-11} hold, in particular, for the Riesz means \eqref{eq:c7-2-14}, when $\alpha$ is larger than the critical index, and for the Abel--Poisson means \eqref{eq:c7-2-13}.
\end{corollary}

The treatment of the Riesz means of order not greater than the critical index is subtler and will be discussed in sections four and five below.

The application of the Poisson summation formula we have just presented concerned periodic analogs of regularization operators arising in the problem of Fourier inversion. The second application is typical of the situation where the Poisson summation formula can be used to give precise estimates involving various elementary (periodic) functions.

In the fourth chapter (see Theorem \ref{thm:c4-4-1}) we learned that if $\alpha$ is a complex number satisfying $0 < \text{Re}(\alpha) < n$ and $f(x) = |x|^{\alpha-n}$, then $\hat{f}(x) = \gamma_\alpha|x|^{-\alpha}$ with $\gamma_\alpha = \pi^{-\alpha+n/2}\Gamma(\alpha/2)/\Gamma[(n-\alpha)/2]$. If we were to apply identity \eqref{eq:c7-2-7}, ignoring all questions of convergence, we would obtain

\begin{equation}\label{eq:c7-2-16}
\gamma_{\alpha}^{-1}\sum_{m\in\Lambda}|x+m|^{\alpha-n}=\sum_{m\in\Lambda}|m|^{-\alpha}e^{2\pi i m\cdot x}.
\end{equation}

This relation, as it stands, cannot be correct because of the infinite term corresponding to $m = 0$ in the right side, and because of the divergence of the series on the left. Nevertheless, the attempt to apply the Poisson summation formula in this context can be made to work after suitable modification. One way of doing this involves the functional equation of the Riemann zeta function and some of its generalizations. This will be described in \eqref{frp:c7-6-3} below. Another interpretation of \eqref{eq:c7-2-16} is given by the following result:

\begin{theorem}
\label{thm:c7-2-17}
Let $0 < \text{Re}(\alpha) < n$. Then the series $\sum_{|m| > 0} |m|^{-\alpha} e^{2\pi i m \cdot x}$ is the Fourier series of an integrable function on $Q_n$ which is of class $C^x$ on $Q_n - \{0\}$. At the origin this function has the same singularity as the function whose values are $\gamma_{\alpha}^{-1}|x|^{\alpha - n}$. That is,
\[
\gamma_{\alpha}^{-1}\vert x\vert^{\alpha-n}+b(x)\sim\sum_{\vert m\vert>0}\vert m\vert^{-\alpha}e^{2\pi i m\cdot x},\qquad(x\in Q_{n})
\]

where $b \in C^{\infty}(Q_n)$.
\end{theorem}

\textsc{Proof.}\ Let us choose a function $\eta$ having the following properties: $\eta \in C^\infty(E_n)$, $\eta(x) = 1$ if $|x| \geq 1$, and $\eta$ vanishes in a neighborhood of the origin. We then define $F$ by letting $F(x) = \eta(x)|x|^{-\alpha}$ for $x \in E_n$. We claim that there exists a function $f \in L^1(E_n)$ such that

(1) $\hat{f} = F$;

(3) For each $n$-tuple $\beta = (\beta_1, \beta_2, \ldots, \beta_n)$ of nonnegative integers and every positive integer $N$, $|D^\beta f(x)| = O(|x|^{-N})$ as $|x| \to \infty$.

In fact, write $F(x) = |x|^{-\alpha} + (\eta(x) - 1)|x|^{-\alpha}$ and let $f$ be the inverse Fourier transform of $F$ in the sense of tempered distributions. Then, by \eqref{thm:c4-4-1} of Chapter IV $f(x) = \gamma_{\alpha}^{-1}|x|^{\alpha-n} + b_1(x)$, where $b_1$ is the inverse Fourier transform of the integrable function with bounded support whose values are $(\eta(x) - 1)|x|^{-\alpha}$. Therefore, $b_1 \in C^\infty(E_n)$ and properties (1) and (2) are established. Since $f$ is the inverse Fourier transform of $F$, for any $n$-tuple $\beta$ of nonnegative integers, $-(2\pi i)^{\beta f}(x)$ is the inverse Fourier transform of $D^\beta F$ (see Theorem \ref{thm:c1-1-8} of Chapter I). Note that if the order of the partial derivative $D^\beta F$ is large enough ( $\beta_1 + \cdots + \beta_n > n - \text{Re}(\alpha)$), then $D^\beta F \in L^1(E_n)$; consequently, when this is the case, $-(2\pi i)^{\beta f}(x)$ is bounded. This shows that $|f(x)| = O(|x|^{-N})$ as $|x| \to \infty$. The rest of property (3) is proved in the same way. The fact that $f \in L^1(E_n)$ follows from (2) and the just established fact that $|f(x)| = O(|x|^{-N})$ as $|x| \to \infty$.

We next apply the Poisson summation formula in the variant given by Theorem \ref{thm:c7-2-4} to the function f we just constructed and $\hat{f} = F$. We obtain

\[
\sum_{m\in\Lambda}f(x+m)\sim\sum_{m\in\Lambda}F(m)e^{2\pi i x\cdot m}=\sum_{|m|>0}|m|^{-\alpha}e^{2\pi i x\cdot m}.
\]

Since

\[
\begin{aligned}\sum_{m\in\Lambda}f(x+m)=f(x)+\sum_{|m|>0}f(x+m)&=\gamma_{\alpha}^{-1}\big|x\big|^{\alpha-n}\\&\quad+b_{1}(x)+\sum_{|m|>0}f(x+m),\end{aligned}
\]

this proves the theorem with $b(x) = b_1(x) + \sum_{|m| > 0} f(x + m)$.

Various other special series besides $\sum_{|m|>0}|m|^{-\alpha}e^{2\pi im\cdot x}$ can be treated by this method. Some of these are described in \eqref{frp:c7-6-1} below.

\section{Multiplier Transformations}

In Chapter I we studied the class of bounded operators from $L^p(E_n)$ into $L^q(E_n)$ which commute with translations. We found that each such operator $T$ is, in effect, determined by a certain tempered distribution $u$ having the property that $Tf = u * f$ for each $f \in \mathcal{L}$. Taking Fourier transforms we then have $(Tf)^\wedge = \hat{u}^f$, so that $T$ is equivalent, via the Fourier transform, to the operator which maps each testing function into the product of the Fourier transform of the testing function with the distribution $\hat{u}$. In general, not too much is known about such ``multipliers'' $\hat{u}$; however, when $p = q$ it can be shown that among other things $\hat{u}$ is a bounded measurable function. Theorem \ref{thm:c1-3-18} of Chapter I asserts that this is the case when $p = 2 = q$. If $p \neq 2$ and $\|u * f\|_p \leq A \|f\|_p$ for all $f \in \mathcal{L}$, it

follows from Theorem \ref{thm:c1-3-20} of Chapter I that $\|u * f\|_p', \leq A \|f\|_p'$ for all $f \in \mathcal{S}$, where $(1/p) + (1/p') = 1$. Since 2 must lie between $p$ and $p'$, it then follows from the M. Riesz interpolation theorem (see Theorem \ref{thm:c5-1-3} of Chapter V) that the operator $T: f \to u \times f$ is of type $(2, 2)$. Consequently $\hat{u}$ satisfies the hypotheses of Theorem \ref{thm:c7-3-18} and, therefore, must be a bounded measurable function.

It is now our intention to study the analogous multiplier operators in the periodic case and to show that for an important class of these their properties are a priori consequences of their nonperiodic variants.

In the notation of Chapter I when $p < \infty$, we let $(L^p, L^q)$ denote the class of bounded operators $T$ from $L^p(E_n)$ into $L^q(E_n)$ which commute with translations (or, equivalently, $(L^p, L^q)$ is the class of tempered distributions $u$ such that $\|u * \varphi\|_q \leq A \|\varphi\|_p$ for some $A = A(u)$ and all $\varphi \in \mathcal{S}$. In order to emphasize the fact that these operators act on functions defined on $E_n$ we shall now write $(L^p(E_n), L^q(E_n))$ for this class. Similarly, for the periodic case, we introduce the class $(L^p(T_n), L^q(T_n))$ of all bounded operators $\tilde{T}$ from $L^p(T_n)$ into $L^q(T_n)$ which commute with translations.

In analogy with the nonperiodic case we could identify any $\tilde{T} \in (L^p(T_n), L^q(T_n))$ with convolution (on $T_n$) by an appropriate tempered distribution. In this case, however, it is much simpler to consider the operators directly.

\begin{theorem}
\label{thm:c7-3-1}
Suppose $\tilde{T} \in (L^p(T_n), L^q(T_n)), 1 \leq p, q \leq \infty$. Then there exists a bounded complex valued function $\lambda: m \to \lambda(m)$ on the lattice $\Lambda$, such that
\begin{equation}\label{eq:c7-3-2}
\tilde{T}f\sim\sum_{m\in\Lambda}\lambda(m)a_{m}e^{2\pi i m\cdot x},
\end{equation}

whenever $f \sim \sum_{m \in \Lambda} a_m e^{2\pi i m \cdot x}$.
\end{theorem}

\textsc{Proof.}\ Set $\psi_m(x) = (\tilde{T} e_m)(x)$, where $e_m(y) = e^{2\pi i m \cdot y}$, for $m \in \Lambda$. Then $\|\psi_m\|_q \leq A \|e_m\|_p = A$. If $\tau_h$ denotes translation by $h$ (that is, for $f$ periodic $(\tau_h f)(x) = f(x - h)$), the hypothesis that $\tilde{T} \tau_h = \tau_h \tilde{T}$ implies that for each $h \in Q_n$

\[
\psi_{m}(x-h)=e^{-2\pi i m\cdot h}\psi_{m}(x)
\]

for almost every $x$. By Fubini's theorem, then, this identity holds for almost every $h \in Q_n$ if a certain $x(= x_0)$ is fixed. This means that $\psi_m(h) = e^{2\pi i m \cdot h} e^{-2\pi i m \cdot x_0} \psi_m(x_0)$ for almost every $h$. Thus,

\[
\psi_{m}(x)=\lambda(m)e^{2\pi i m\cdot x}
\]

for almost every $x$ with $\lambda(m) = e^{-2\pi im \cdot x_0} \psi_m(x_0)$. Moreover, $|\lambda(m)| = |\lambda(m)|\|e_m\|_q = \|\lambda(m)e_m\|_q = \|\psi_m\|_q \leq A$. Passing to finite linear combinations of the functions $\{e_m\}$ we obtain the representation \eqref{eq:c7-3-2} for $f$ a trigonometric polynomial. By a simple limiting argument this extends to all $f \in L^p(T_n)$ and the theorem is established.

Because of Theorem \ref{thm:c7-3-1} we see that it is natural to refer to the transformations $\widetilde{T} \in (L^p(T_n), L^q(T_n))$ as multiplier\index{multiplier operators@multiplier operators!associated with n-torus@associated with $n$-torus} operators and the associated sequences $\{\lambda(m)\}$ as multipliers\index{multipliers@multipliers}.

The following corollary is a consequence of this theorem and part (iii) of \eqref{thm:c7-1-7}:

\begin{corollary}
\label{cor:c7-3-3}
$\tilde{T}\in(L^2(T_n),L^2(T_n))$ if and only if $\{\lambda(m)\}$ is a bounded sequence. Moreover, the operator norm $\|\tilde{T}\|$ equals $\sup_{m\in\Lambda}|\lambda(m)|$.
\end{corollary}

The $L^{1}$ case is similar to the nonperiodic one treated in the first chapter. More precisely, we show the following result:

\begin{theorem}
\label{thm:c7-3-4}
$\tilde{T} \in (L^1(T_n), L^1(T_n))$ if and only if there exists a finite Borel measure $\mu$ on $T_n$, such that
\begin{equation}\label{eq:c7-3-5}
d\mu\sim\sum_{m\in\Delta}\lambda(m)e^{2\pi i m\cdot x}.
\end{equation}

\end{theorem}

In this case $\|d_\mu\| = \|\tilde{T}\|$.

\textsc{Proof.}\ Suppose $\|\tilde{T}f\|_1 \leq A \|f\|_1$ for all $f \in L^1(T_n)$. For $t > 0$ let $f_t(x) = P_t(x) = \sum_{m \in \Lambda} e^{-2\pi |m| t} e^{2\pi i m \cdot x}$. We showed earlier (see the discussion following \eqref{eq:c7-2-13}) that $\|f_t\|_1 = 1$. Thus,

\[
\|\tilde{T}f_{t}\|_{1}=\int\limits_{Q_{n}}\left|\sum_{m\in\Lambda}e^{-2\pi|m|t}e^{2\pi i m\cdot x}\lambda(m)\right|d x\leq A.
\]

Consequently, we can find a sequence $\{t_n\}$ of positive numbers tending to 0 such that $\{\tilde{T}f_{t_n}\}$ tends weakly to a measure $\mu$. That is,

\[
\lim_{j\to\infty}\int\limits_{T_{n}}\tilde{T}f_{t_{j}}(x)g(x)d x=\int\limits_{T_{n}}g(x)d\mu(x)
\]

for each $g \in C(T_n)$. It follows that $\|d\mu\| \leq A$. In particular,

\[
\lambda(m)=\lim_{j\to\infty}\int_{T_{n}}\mathopen{e}^{-2\pi i m\cdot x}\widetilde{T}f_{t_{j}}(x)\mathop{}{d}x=\int_{T_{n}}\mathopen{}{e}^{-2\pi i m\cdot x}\mathop{}{d}\mu(x).
\]

Therefore,

\[
d\mu\sim\sum_{m\in\Lambda}\lambda(m)e^{2\pi i m\cdot x}.
\]

The converse is an immediate consequence of the definition of the convolution operator (see \eqref{eq:c7-1-6} and the discussion that follows).

Having established these preliminary results we turn to the main question to be discussed in this section. Suppose $T$ belongs to $(L^p(E_n), L^p(E_n))$. We noticed above that $T$ could be realized in terms of a multiplier

$\hat{u}$, which is a bounded measurable function. In general, $\hat{u}$ is not continuous and so $\hat{u}$ need not be defined on the points of the lattice $\Lambda$. Suppose, however, that $\hat{u}$ were continuous at the points of $\Lambda$, so that $\lambda(m) = \hat{u}(m)$ is well-defined on $\Lambda$. We can then ask whether $\{\lambda(m)\}$ is a multiplier in the class $(L^p(T_n), L^p(T_n))$. If this is the case we shall say that the corresponding operator $\tilde{T}$ on the $n$-torus arises from the operator $T$ by periodization\index{periodization@periodization!of operators@of operators}.

Before we formulate the general theorem we examine some of the simpler special cases. When $p = 2$, the conditions that $\hat{u}$ is bounded on $E_n$ and continuous at the points of $\Lambda$ imply $\sup_{m \in \Lambda} |\hat{u}(m)| = \sup_{m \in \Lambda} |\lambda(m)| < \infty$. By Corollary \ref{cor:c7-3-3}, therefore, the periodized operator belongs to $(L^2(T_n), L^2(T_n))$. For the case $p = 1$, we know that $\hat{u}(x) = \hat{\mu}(x)$, where $\hat{\mu}$ is the Fourier transform of a measure $\mu \in \mathcal{B}(E_n)$ (see Theorem \ref{thm:c1-3-19} of Chapter I). The fact that the corresponding periodized operator belongs to $(L^1(T_n), L^1(T_n))$ is then a consequence of Theorem \eqref{thm:c7-3-4} and the following observation (which can be considered as still another variant of the Poisson summation formula):

\begin{theorem}
\label{thm:c7-3-6}
Suppose $\mu \in \mathcal{B}(E_n)$ and $\hat{\mu}$ is its Fourier transform. Then $\sum_{m \in \Lambda} \hat{\mu}(m) e^{2\pi i m \cdot x}$ is the Fourier series of a measure $\tilde{\mu}$ on $T_n$; moreover, $\|d\tilde{\mu}\| \leq \|d\mu\|$.
\end{theorem}

\textsc{Proof.}\ Consider the linear functional mapping $f \in C(T_n)$ into $\int_{E_n} f(x) \, d\mu(x)$ (we are using the symbol $f$ to denote the periodic extension of $f$ as well). This functional can be realized, by the Riesz representation theorem, in terms of some measure $\tilde{\mu} \in \mathcal{B}(T_n)$. That is,

\begin{equation}\label{eq:c7-3-7}
\int\limits_{E_{n}}f(x)d\mu(x)=\int\limits_{T_{n}}f(x)d\widetilde{\mu}(x).
\end{equation}

Theorem \ref{thm:c7-3-6} now follows by applying \eqref{eq:c7-3-7} to $f(x) = e^{-2\pi im\cdot x}$ for $m \in \Lambda$.

Observe that it follows from \eqref{eq:c7-3-7} that for any Borel set $E \subset Q_n$, $\tilde{\mu}(E) = \sum_{m \in \Lambda} \mu(E + m)$. Thus, the measure $\tilde{\mu}$ arises from $\mu$ by a method of periodization analogous to that we introduced for functions (see \eqref{eq:c7-2-1} and the discussion following).

The main result of this section can be formulated as follows:

\begin{theorem}
\label{thm:c7-3-8}
Suppose $1 \leq p \leq \infty$ and $T \in (L^p(E_n), L^p(E_n))$. Let $\hat{u}$ be the multiplier corresponding to $T$ and suppose that $\hat{u}$ is continuous at each point of the lattice $\Lambda$. Set $\lambda(m) = \hat{u}(m)$ for $m \in \Lambda$. Then there exists a unique periodized operator $\tilde{T}$ defined by \eqref{eq:c7-3-2} such that $\tilde{T} \in (L^p(T_n), L^p(T_n))$ and $\|\tilde{T}\| \leq \|T\|$.
\end{theorem}

We begin the proof by considering two lemmas.

\begin{lemma}
\label{lem:c7-3-9}
Suppose f is a continuous periodic function on $E_{n}$, then
\begin{equation}\label{eq:c7-3-10}
\lim_{\varepsilon\to0}\;\varepsilon^{n/2}\int_{E_{n}}f(x)e^{-\varepsilon\pi|x|^{2}}d x=\int_{Q_{n}}f(x)d x.
\end{equation}

\end{lemma}

\textsc{Proof.}\ Equality \eqref{eq:c7-3-10} is obvious when $f(x) = e^{2\pi im \cdot x}$ since (see \eqref{thm:c1-1-13} of Chapter I)

\[
\varepsilon^{n/2}\int\limits_{E_{n}}e^{2\pi i m\cdot x}e^{-\varepsilon\pi|x|^{2}}d x=e^{-\pi|m|^{2}/\varepsilon}
\]

for $\varepsilon > 0$ and $m \in \Lambda$. Hence, \eqref{eq:c7-3-10} holds for any trigonometric polynomial. The lemma now follows by approximating an arbitrary continuous periodic function uniformly on $Q_n$ by such polynomials. The second lemma, which is the principal tool we shall use in the proof of Theorem \ref{thm:c7-3-8}, is the following:

\begin{lemma}
\label{lem:c7-3-11}
Suppose $P$ and $Q$ are trigonometric polynomials, $T$ belongs to $(L^p(E_n), L^p(E_n))$, $\tilde{T}$ is defined on the class of trigonometric polynomials by \eqref{eq:c7-3-2} and $w_\delta(y) = e^{-\pi\delta|y|^2}$ for $\delta > 0$, $y \in E_n$. Then
\begin{equation}\label{eq:c7-3-12}
\lim_{\varepsilon\to0}\;\varepsilon^{n/2}\int_{E_{n}}T(P w_{\varepsilon\alpha})(x)\overline{{Q(x)}}w_{\varepsilon\beta}(x)\;d x=\int_{Q_{n}}(\tilde{T}P)(x)\overline{{Q(x)}}\;d x,
\end{equation}

whenever $\alpha, \beta > 0$ and $\alpha + \beta = 1$.
\end{lemma}

\textsc{Proof.}\ Since the expressions in \eqref{eq:c7-3-12} are linear in $P$ and $Q$ it suffices to prove \eqref{eq:c7-3-12} when $P(x) = e^{2\pi i m \cdot x}$ and $Q(x) = e^{2\pi i k \cdot x}$ for $m, k \in \Lambda$. By Plancherel's theorem and the definition of the multiplier $\hat{u}$, the integral in the left side equals $\varepsilon^{n/2} \int_{E_n} \hat{u}(x) \varphi(x) \overline{\psi(x)} \, dx$, when $\varphi$ and $\psi$ are the Fourier transforms of the functions having values $e^{2\pi i m \cdot x} e^{-\pi\varepsilon\alpha |x|^2}$ and $e^{2\pi i k \cdot x} e^{-\pi\varepsilon\beta |x|^2}$. But, by \eqref{eq:c1-1-5} and Theorem \ref{thm:c1-1-13} of Chapter I we have

\[
\varphi(x)=e^{-\pi(|x-m|^{2}/\alpha\varepsilon)}(\alpha\varepsilon)^{-(n/2)}\quad\text{and}\quad\psi(x)=e^{-\pi(|x-k|^{2}/\beta\varepsilon)}(\beta\varepsilon)^{-(n/2)}.
\]

Assume now that $m \neq k$ and, consequently, $|m - k| \geq 1$. Since $|\hat{u}(x)| \leq A$ for a suitable constant $A$, the left side of \eqref{eq:c7-3-12} is bounded by

\[
\begin{aligned}\varepsilon^{n/2}A\int\limits_{E_{n}}e^{-(|x-m|^{2}/\alpha\varepsilon)\pi}\big(\alpha\varepsilon\big)^{-(n/2)}&e^{-(|x-k|^{2}/\beta\varepsilon)\pi}\big(\beta\varepsilon\big)^{-(n/2)}dx\\\ &\leq\varepsilon^{n/2}A\Bigg[\int\limits_{|x-m|\geq1/2}+\int\limits_{|x-k|\geq1/2}\Bigg].\end{aligned}
\]

In the integral extended over $\{x \in E_n: |x - m| \geq \frac{1}{2}\}$, the factor $\varepsilon^{n/2} e^{-(|x - m|^2/\alpha\varepsilon)\pi}(\alpha\varepsilon)^{-(n/2)}$ tends uniformly to 0 as $\varepsilon$ tends to 0, while

the factor $e^{-(|x-k|^{\varepsilon}/\beta\varepsilon)\pi}(\beta\varepsilon)^{-(n/2)}$ has total integral 1 when extended over $E_n$. It follows, therefore, that $\varepsilon^{(n/2)}\int_{|x-m|\geq1/2}$ tends to 0 with $\varepsilon$. The same argument, with the roles of $m$ and $k$ interchanged shows that $\lim_{\varepsilon\to0}\varepsilon^{n/2}\int_{|x-k|\geq1/2}=0$. Since, for $m\neq k$,

\[
\int_{T_{n}}(\widetilde{T}P)(x)\overline{{Q(x)}}d x=\int_{Q_{n}}\lambda(m)e^{2\pi\mathsf{i}m\cdot x}e^{-2\pi\mathsf{i}k\cdot x}d x=0,
\]

equality \eqref{eq:c7-3-12} is established for this case.

In case m = k, the left side of \eqref{eq:c7-3-12} equals

\begin{equation}\label{eq:c7-3-13}
\lim_{\varepsilon\to0}\;\big(\varepsilon\alpha\beta\big)^{-(n/2)}\;\int\limits_{E_{n}}\hat{u}(x)e^{-\pi(|x-m|^{2}/\varepsilon)(1/\alpha+1/\beta)}\;d x.
\end{equation}

Since $1/\alpha + 1/\beta = 1/\alpha\beta$, \eqref{eq:c7-3-13} is the limit, as $\varepsilon \to 0$, of the Gauss-Weier-strass integral of $\hat{u}$. By Theorem \ref{thm:c1-1-25} of Chapter I, therefore, this limit is $\hat{u}(m)$ provided $m$ belongs to the Lebesgue set of $\hat{u}$. But this is the case since $\hat{u}$ is assumed to be continuous at $m$. This proves equality \eqref{eq:c7-3-12} when $P(x) = e^{2\pi im \cdot x} = Q(x)$ (since, in this case, $\int_{T_n}(\tilde{T}P)(x)\overline{Q(x)} \, dx = \lambda(m)$) and the lemma is established.

We now pass to the proof of Theorem \ref{thm:c7-3-8}. In order to avoid certain irrelevant technicalities we shall temporarily assume that $1 < p < \infty$. Let $q$ be the conjugate exponent to $p$; then $1/p + 1/q = 1$ and $1 < q < \infty$. We first prove that there exists $A \leq \|T\|$ such that

\begin{equation}\label{eq:c7-3-14}
\left(\intop_{Q_{n}}\left|(\tilde{T}P)(x)\right|^{p}d x\right)^{1/p}\leq A\left(\intop_{Q_{n}}\left|P(x)\right|^{p}d x\right)^{1/p}
\end{equation}

for all trigonometric polynomials P. If Q is also a trigonometric polynomial then

\begin{equation}\label{eq:c7-3-15}
\left|\int_{E_{n}}{(T(P w_{\varepsilon\alpha}))(x)\overline{{Q(x)}}w_{\varepsilon\beta}(x)\;}d x\right|\;\leq\;\left\|T\right\|\;\left\|P w_{\varepsilon\alpha}\right\|_{p}\left\|Q w_{\varepsilon\beta}\right\|_{q},
\end{equation}

where the norms are taken with respect to $E_n$ and $w_\delta$, for $\delta > 0$, is the function introduced in Lemma \ref{lem:c7-3-11}. Set $\alpha = (1/p)$, $\beta = (1/q)$ multiply both sides by $\varepsilon^{n/2}$ and let $\varepsilon \to 0$. By Lemma \ref{lem:c7-3-11} the left side converges to $\int_{Q_n}(\tilde{T}P)(x)\overline{Q(x)}$ dx. On the other hand, by Lemma \ref{lem:c7-3-9}

\[
\begin{aligned}&\lim_{\varepsilon\to0}\varepsilon^{n/2}\big\|P w_{\varepsilon/p}\big\|_{p}\big\|Q w_{\varepsilon/q}\big\|_{q}\\&=\lim_{\varepsilon\to0}\left[\varepsilon^{n/2}\int\limits_{E_{n}}\big|P(x)\big|^{p} e^{-\varepsilon\pi|x|^{2}}dx\right]^{1/p}\left[\varepsilon^{n/2}\int\limits_{E_{n}}\big|Q(x)\big|^{q} e^{-\varepsilon\pi|x|^{2}}dx\right]^{1/q}\\&=\left[\int\limits_{Q_{n}}\big|P(x)\big|^{p} dx\right]^{1/p}\left[\int\limits_{Q_{n}}\big|Q(x)\big|^{q} dx\right]^{1/q}.\end{aligned}
\]

Combining this with \eqref{eq:c7-3-15} we have

\[
\bigg|\int_{Q_{n}}(\widetilde T P)(x)\overline{{Q(x)}}\;d x\bigg|\;\leq\;\Vert T\Vert\left(\int_{Q_{n}}\vert P(x)\vert^{p}\;d x\right)^{1/p}\left(\int_{Q_{n}}\vert Q(x)\vert^{q}\;d x\right)^{1/q}.
\]

Finally, taking the supremum over all polynomials $Q$ satisfying $\int_{Q_n}|Q(x)|^q \, dx \leq 1$, we obtain \eqref{eq:c7-3-14}. This shows that the restriction of $\tilde{T}$ to trigonometric polynomials is a bounded operator on $L^p(T_n)$ with bound not exceeding $\|T\|$. This restriction then has a unique bounded extension to all of $L^p(T_n)$, and it is this extension that satisfies the conclusions of Theorem \ref{thm:c7-3-8}.

We now pass to the cases $p = 1$ and $p = \infty$. These extreme cases are, in fact, much simpler than the general one we just considered. The result for $p = 1$ is (as we have already pointed out) an immediate consequence of Theorem \ref{thm:c1-3-19} of Chapter I and Theorems \ref{thm:c7-3-4} and 3.6 of this chapter. When $p = \infty$ we argue as follows: By Theorem \ref{thm:c1-3-20} of Chapter I, if $T \in (L^\infty(E_n), L^\infty(E_n))$, then $T \in (L^1(E_n), L^1(E_n))$; moreover, the proof of Theorem \ref{eq:c7-3-20} shows that the norm of $T$, as an operator on $L^1(E_n)$, is no larger than its norm as an operator on $L^\infty(E_n)$. Thus from the just established case $p = 1$, $\widetilde{T} \in (L^1(T_n), L^1(T_n))$ and, by \eqref{thm:c7-3-4}, there exists a finite Borel measure $\mu$ on $T_n$ such that $\|d_\mu\| = \|\widetilde{T}\|$ and $\widetilde{T}f = f^* d_\mu$. For $f \in L^\infty(T_n)$, however, $\|f^* d_\mu\|_\infty \leq \|f\|_\infty$ and $d_\mu\|$ and the theorem is also proved in this case.

\begin{corollary}
\label{cor:c7-3-16}
The conclusions of Theorem \ref{thm:c7-3-8} are still valid if, instead of assuming $\hat{u}$ is continuous at the points of $\Lambda$, we assume that
\begin{equation}\label{eq:c7-3-17}
\lim_{\varepsilon\to0}\varepsilon^{-n}\int\limits_{|t|\leq\varepsilon}\left[\hat{u}(m-t)-\hat{u}(m)\right]d t=0
\end{equation}

for each $m \in \Lambda$.
\end{corollary}

In fact, the assumption \eqref{eq:c7-3-17} suffices to show that the limit \eqref{thm:c1-3-13} exists and has value $\hat{u}(m)$ (see \eqref{eq:c1-4-10} in Chapter I).

An important class of these multiplier transforms arises from the singular integral operators that were studied in the last chapter. There we showed (see also Theorem \ref{thm:c4-4-7}, Corollary \ref{cor:c4-4-12}, and \eqref{frp:c4-5-10} of Chapter IV) that the multiplier corresponding to a singular integral operator $T$ is a homogeneous function of degree 0, $\Omega_0$, which is continuous on the surface of the unit sphere (hence, continuous everywhere but the origin) and, moveover, has the property $\int_{\Sigma_{n-1}} \Omega_0(x') \, dx' = 0$. If we set $\hat{u}(x) = \Omega_0(x)$, for $x \neq 0$, and $\hat{u}(0) = 0$, then condition \eqref{eq:c7-3-17} of Corollary \ref{cor:c7-3-16} is

satisfied; consequently, the operator $\tilde{T}$ defined by

\[
(T f)(x)\sim\sum_{m\neq0}\Omega_{0}(m)a_{m}e^{2\pi i m\cdot x}
\]

whenever

\[
f(x)\sim\sum_{m\in\Delta}a_{m}e^{2\pi i m\cdot x}
\]

belongs to $(L^p(T_n), L^p(T_n))$ for $1 < p < \infty$. A particular example arises when $\Omega_0(x) = P^{(k)}(x)/|x|^k$, where $P^{(k)}$ is a harmonic polynomial on $E_n$ that is homogeneous of degree $k \geq 1$ (see Theorem \ref{thm:c4-4-5} of Chapter IV). The special cases $\Omega_0(x) = -ix_j/|x|$, $j = 1, 2, \ldots, n$, are known as the periodic Riesz transforms\index{Riesz transforms@Riesz transforms!periodic@periodic} (see \eqref{eq:c6-2-7} of Chapter VI).

We shall now show that Theorem \ref{thm:c7-3-8} has a converse. Suppose $\lambda$ is a continuous function on $E_n$ and the family $\{\lambda(m)\}, m \in \Lambda$, consists of the multipliers of an operator in $(L^p(T_n), L^p(T_n))$. Can we conclude that $\lambda$ is a multiplier of an operator in $(L^p(E_n), L^p(E_n))$? It is easy to see that this cannot be the case since the assumption we are making on $\lambda$ bears only on its values at the points of the lattice $\Lambda$, while the conclusion would involve its behavior on all of $E_n$. To formulate a meaningful converse we observe that if $\lambda$ were a multiplier of an operator in $(L^p(E_n), L^p(E_n))$ then, for each $\varepsilon > 0$, the function $\lambda(\varepsilon x)$ would also be a multiplier of an operator $T$ of type $(L^p(E_n), L^p(E_n))$ whose norm depends only on $\lambda$ and not on $\varepsilon$. In fact, we have $(Tf)^\gamma(x) = \lambda(x) \hat{f}(x)$ for each $f \in \mathcal{S}$. Thus, if we define $T_\varepsilon$ to be the operator $\delta_\varepsilon^{-1}T$ $\delta_\varepsilon$ ( $\delta_\varepsilon$ being the dilation operator introduced immediately before \eqref{eq:c7-1-6} of Chapter 1) we obtain again a multiplier operator whose multiplier is the function having the values $\lambda(\varepsilon x)$. Moreover, since $\|\delta_\varepsilon f\|_p = \varepsilon^{-n/p} \|f\|_p$ and $\|\delta_\varepsilon^{-1}f\|_p = \varepsilon^{n/p} \|f\|_p$, we see that $\|T_\varepsilon\| = \|T\|$. In view of this observation, it is not surprising that the following is a converse of \eqref{thm:c7-3-8}:

\begin{theorem}
\label{thm:c7-3-18}
Let $\lambda$ be a continuous function on $E_n$. Suppose that for each $\varepsilon > 0$ there exists an operator $\tilde{T}_\varepsilon \in (L^p(T_n), L^p(T_n))$ given by
\begin{equation}\label{eq:c7-3-19}
(T_{\varepsilon}f)(x)\sim\sum_{m\in\Lambda}\lambda(\varepsilon m)a_{m}e^{2\pi i m\cdot x}
\end{equation}

whenever $\{a_m\}$ are the Fourier coefficients of $f \in L^p(T_n)$. We assume that the norms $\|\tilde{T}_\varepsilon\|$ of the operators $\tilde{T}_\varepsilon$ are uniformly bounded. Then $\lambda$ is a multiplier of type $(L^p(E_n), L^p(E_n))$; moreover, if $T$ is the corresponding operator, its operator norm $\|T\|$ does not exceed $\sup \|\tilde{T}_\varepsilon\|$.
\end{theorem}

\textsc{Proof.}\ We dispose first of the case $p = \infty$ by showing that it can be reduced to the case $p = 1$. In fact, \eqref{eq:c7-3-19} implies immediately the duality

identity

\begin{equation}\label{eq:c7-3-20}
\intop_{T_{n}}(\tilde{T}_{\varepsilon}f)(x)g(-x)d x=\intop_{T_{n}}(\tilde{T}_{\varepsilon}g)(x)f(-x)d x
\end{equation}

whenever $f$ and $g$ are (say) trigonometric polynomials. From this it follows (see the analogous argument in the proof of Theorem \ref{eq:c7-3-20} of Chapter 1) that $\|\widetilde{T}_{\varepsilon}\|_{1} \leq \|\widetilde{T}_{\varepsilon}\|_{\infty}$ (the subscripts are used here to denote the norms of $\widetilde{T}_{\varepsilon}$ as an operator on $L^{1}(T_{n})$ and on $L^{\infty}(T_{n})$, respectively). Thus, if the case $p=1$ of this theorem were established, we would obtain an operator $T \in (L^{1}(E_{n}), L^{1}(E_{n}))$, having multiplier $\lambda$, satisfying $\|T\|_{1} \leq \sup_{\varepsilon > 0} \|\widetilde{T}_{\varepsilon}\|_{1} \leq \sup_{\varepsilon > 0} \|\widetilde{T}_{\varepsilon}\|_{\infty}$. Finally, again appealing to Theorem \ref{thm:c1-3-20} of Chapter I, we would obtain the fact that $T \in (L^{\infty}(E_{n}), L^{\infty}(E_{n}))$ and $\|T\|_{\infty} \leq \sup_{\varepsilon > 0} \|\widetilde{T}_{\varepsilon}\|_{\infty}$.

We shall, therefore, assume that $1 \leq p < \infty$ and make use of the following appropriate partition of unity:

\begin{lemma}
\label{lem:c7-3-21}
There exists a nonnegative continuous function $\eta$ having compact support in $E_n$ and satisfying
(a) $\eta(0) = 1$;
\begin{equation}
\tag{b}
\textstyle\sum_{m\in\Lambda}\;[\eta(x+m)]^{p}\equiv1.
\end{equation}

\end{lemma}

To prove this lemma choose any nonnegative continuous function $\eta_1$ with compact support in $E_n$ such that $\eta_1(0) = 1$, $\eta_1(m) = 0$ if $m \in \Lambda - \{0\}$ and $\eta_1(x) > 0$ for $x \in \widetilde{Q}_n$. Let $\eta_2(x) = \eta_1(x)/\sum_{m \in \Lambda} \eta_1(x + m)$. Then, clearly, $\eta_2(0) = 1$ and $\sum_{m \in \Lambda} \eta_2(x + m) \equiv 1$. We now need only take $\eta = \eta_2^{1/p}$.

Let us now return to the proof of the theorem. We assume, for simplicity, that $\|\widetilde{T}_\varepsilon\|_p \leq 1, \varepsilon > 0$. Then $|\lambda(\varepsilon m)| \leq 1$ for $m \in \Lambda$ and $\varepsilon > 0$ (see the proof of Theorem \ref{thm:c7-3-1}). Since the set $\{\varepsilon m\colon \varepsilon > 0, m \in \Lambda\}$ is dense in $E_n$ it follows that $\lambda$ is bounded. When $f \in L^2(E_n)$, therefore, $\lambda^f$ also belongs to $L^2(E_n)$; hence, the latter function is the Fourier transform of a square integrable function. In particular, this allows us to define $Tf$, for $f \in \mathcal{D} (C \setminus S)$, to be the function whose Fourier transform is $\lambda f$; that is, $(Tf)^\cap(x) = \lambda(x)f(x)$. We shall show that

\begin{equation}\label{eq:c7-3-22}
\left\|T f\right\|_{p}\leq\left\|f\right\|_{p}.
\end{equation}

In order to do so, define $\tilde{f}_{\varepsilon}$, for $\varepsilon > 0$, to be the dilated and periodized version of f, viz.,

\[
\tilde{f}_{\varepsilon}(x)=\varepsilon^{-n}\sum_{m\in\Lambda}f\left(\frac{x+m}{\varepsilon}\right).
\]

Then, using the Poisson summation formula \eqref{eq:c7-2-7}, we obtain

\begin{equation}\label{eq:c7-3-23}
\tilde{f}_{\varepsilon}(x)=\sum_{m\in\Lambda}\hat{f}(\varepsilon m)e^{2\pi i m\cdot x}.
\end{equation}

We now claim that, for each $x \in E_n$,

\begin{equation}\label{eq:c7-3-24}
\lim_{\varepsilon\to0}\varepsilon^{n}[\widetilde{T}_{\varepsilon}\widetilde{f}_{\varepsilon}](\varepsilon x)=[T f](x).
\end{equation}

From \eqref{eq:c7-3-19} and \eqref{eq:c7-3-23} we obtain

\begin{equation}\label{eq:c7-3-25}
\varepsilon^{n}[\widetilde{T}_{\varepsilon}\widetilde{f}_{\varepsilon}](\varepsilon x)\;=\;\varepsilon^{n}\sum_{m\in\Lambda}\;\lambda(\varepsilon m)\widehat{f}(\varepsilon m)e^{2\pi i\varepsilon m\cdot x}.
\end{equation}

Now, $\lambda$ is bounded, $\hat{f}$ is rapidly decreasing at $\infty$ (i.e., $\lim_{|x|\to\infty}|x|^{k}|\hat{f}(x)|=0$ for all positive integers $k$) and both $\lambda$ and $\hat{f}$ are continuous. Consequently, by the definition of the Riemann integral, the right side of \eqref{eq:c7-3-25} tends to

\[
\int\limits_{E_{n}}\lambda(t)\hat{f}(t)e^{2\pi\mathrm{i}x\cdot t}dt=(Tf)(x)
\]

as $\varepsilon$ tends to 0 (see Corollary \ref{cor:c1-1-21} of Chapter I). Therefore, equality \eqref{eq:c7-3-24} is established. Moreover, since $\eta$ is continuous and $\eta(0) = 1$, we also have

\begin{equation}\label{eq:c7-3-26}
\lim_{\varepsilon\to0}\varepsilon^{n}[\widetilde{T}_{\varepsilon}\widetilde{f}_{\varepsilon}](\varepsilon x)\gamma(\varepsilon x)=(T f)(x)
\end{equation}

for each $x \in E_n$.

Because of the periodicity of $\tilde{T}_{\varepsilon}\tilde{f}_{\varepsilon}$ we have

\[
\begin{aligned}{\varepsilon^{n p}\int_{E_{n}}|(\widetilde{T}_{\varepsilon}\widetilde{f}_{\varepsilon})(\varepsilon x)\eta(\varepsilon x)|^{p}d x}&{{}=\varepsilon^np-n\int_{E_{n}}|(\widetilde{T}_{\varepsilon}\widetilde{f}_{\varepsilon})(x)|^{p}[\eta(x)]^{p}d x}\\ {}&{{}=\varepsilon^np-n\sum_{m\in\Lambda}\int_{Q_{n}}|(\widetilde{T}_{\varepsilon}\widetilde{f}_{\varepsilon})(x)|^{p}[\eta(x+m)]^{p}d x.}\\ \end{aligned}
\]

Thus, by Lemma \ref{lem:c7-3-21} and the assumption $\|\tilde{T}_{\varepsilon}\|_{p} \leq 1$,

\begin{equation}\label{eq:c7-3-27}
\int_{{E}_{n}}|\varepsilon^{n}(\widetilde{T}_{\varepsilon}\widetilde{f}_{\varepsilon})(\varepsilon x)\eta(\varepsilon x)|^{p}\;{d}x\;\leq\;\varepsilon^np-n\int_{{Q}_{n}}|\widetilde{f}_{\varepsilon}(x)|^{p}\;{d}x.
\end{equation}

For $\varepsilon$ sufficiently small the support of $\varepsilon^{-n}f(x/\varepsilon)$ lies entirely in $Q_n$ and, in this case, $\varepsilon^{-n}f(x/\varepsilon) = \tilde{f}_{\varepsilon}(x)$ for $x \in Q_n$. Thus, for small $\varepsilon$ the right side of \eqref{eq:c7-3-27} equals

\[
\varepsilon^np-n\intop_{Q_{n}}|\varepsilon^{-n}f(x/\varepsilon)|^{p}d x=\varepsilon^np-n\intop_{E_{n}}|\varepsilon^{-n}f(x/\varepsilon)|^{p}d x=\intop_{E_{n}}|f(x)|^{p}d x.
\]

Now, applying \eqref{eq:c7-3-26} and Fatou's lemma to the left side of \eqref{eq:c7-3-27} we obtain

\[
\int\limits_{E_{n}}\left|(T f)(x)\right|^{p}d x\leq\int\limits_{E_{n}}\left|f(x)\right|^{p}d x,
\]

which is the desired inequality \eqref{eq:c7-3-22}. Since the class $\mathcal{D}$ is dense in $L^p(E_n)$, $1 \leq p < \infty$, this proves the theorem.

Theorem \ref{thm:c7-3-18}, the remarks immediately preceding it, and Theorem \ref{thm:c7-3-8} have the following obvious corollary:

\begin{corollary}
\label{cor:c7-3-28}
$\sup_{\varepsilon>0} \|\tilde{T}_\varepsilon\| = \|T\|$, where $\|\tilde{T}_\varepsilon\|$ and $\|T\|$ denote the operator norms of the operators $\tilde{T}_\varepsilon$ and $T$ occurring in Theorem \ref{thm:c7-3-18}.
\end{corollary}

\section{Summability Below the Critical Index (Negative Results)}

Let $\sum_{m \in \Lambda} a_m e^{2\pi i m \cdot x}$ be the Fourier series of an integrable function $f$. COROLLARY \eqref{cor:c7-2-15} asserts that the equality

\begin{equation}\label{eq:c7-4-1}
\lim_{R\to\infty}\sum_{|m|<R}\left(1-\frac{|m|^{2}}{R^{2}}\right)^\alpha a_{m} e^{2\pi i m\cdot x}=f(x)
\end{equation}

is valid (almost everywhere and also if the limit is taken in the $L^1$ norm), provided $\alpha$ is larger than the critical index $(n - 1)/2$. It is natural to ask what happens when $\alpha \leq (n - 1)/2$. In this section we shall be concerned mainly with the two cases $\alpha = (n - 1)/2$ and $\alpha = 0$.

In order to gain some perspective about this question we begin by reviewing the classical results when $n = 1$. In this case summability at the critical index $(\alpha = 0)$ and ordinary convergence coincide, Kolmogorov showed that there exists an $L^1$ function for which the limit \eqref{eq:c7-4-1} fails to exist for every $x$ when $\alpha = 0$. Nevertheless, there are localization results for $L^1$ functions which assert that \eqref{eq:c7-4-1} holds for a given $x$ as long as $f$ is sufficiently regular in an arbitrarily small neighborhood of $x$. For a function $f$ in $L^p(T_1)$, $p > 1$, there is a result of Carleson and Hunt which shows that the Fourier series of $f$ converges to $f(x)$ almost everywhere. There is also an older result of M. Riesz which shows that convergence in the norm holds when $1 < p < \infty$.

We shall see now that if n > 1, the situation has a rather different character. This is clear from the following three assertions:

\begin{theorem}
\label{thm:c7-4-2}
There exists an $f\in L^{1}(T_{n})$, $n>1$, such that
\[
\lim_{R\to\infty}\sup_{|m|<R}\biggl(1\;-\;\frac{|m|^{2}}{R^{2}}\biggr)^{(n-1)/2}a_{m}e^{2\pi i m\cdot x}\biggr\vert=\infty
\]

for almost every x. Such a function f can be constructed so that it is supported in an arbitrarily small given neighborhood of the origin. \footnote{This theorem shows at once that neither almost everywhere summability nor localization can hold at the critical index.}
\end{theorem}

\begin{theorem}
\label{thm:c7-4-3}
The trigonometric series
\begin{equation}\label{eq:c7-4-4}
\sum_{m\neq0}\left|m\right|^{-(n/2)+1/2}e^{2\pi i m\cdot x}
\end{equation}

diverges almost everywhere. More particularly,
\[
\limsup_{R\to\infty}\bigg\vert\sum_{\substack{0<|m|<R}}|m|^{-(n/2)+1/2}e^{2\pi i m\cdot x}\bigg\vert=\infty
\]

for almost every $x \in Q_n$.
\end{theorem}

We know by Theorem \ref{thm:c7-2-17} that the series \eqref{eq:c7-4-4} is the Fourier series of the function whose values are $\gamma_{(n-1)/2}^{-1}|x|^{-(n+1)/2} + b(x)$, where $b \in C^{\infty}(Q_n)$. Thus, we also have the following corollary:

\begin{corollary}
\label{cor:c7-4-5}
There exists a function which belongs to $L^{p}(T_{n})$, $p < 2n/(n + 1)$, whose Fourier series diverges almost everywhere.
\end{corollary}

Therefore, for the general function of $L^{p}(T_{n}), p < 2$, convergence almost everywhere fails if the dimension $n$ is sufficiently large. It is an open problem to determine what happens for $p = 2$.

These three results are based on the following lemma:

\begin{lemma}
\label{lem:c7-4-6}
If n > 1 then
\begin{equation}\label{eq:c7-4-7}
\limsup_{R\to\infty}\left|\sum_{|m|<R}\left(1-\frac{|m|^{2}}{R^{2}}\right)^{(n-1)/2}e^{2\pi im\cdot x}\right|=\infty
\end{equation}

for almost every $x \in Q_{n}$.
\end{lemma}

We break up the proof of \eqref{lem:c7-4-6} into several steps. In doing so we set

\[
K_{R}^{\alpha}(x)=\sum_{|m|<R}\left(1\;-\;\frac{|m|^{2}}{R^{2}}\right)^{\alpha}e^{2\pi i m\cdot x}
\]

and shall often write $K_{R}$ for $K_{R}^{(n-1)/2}$.

\begin{lemma}
\label{lem:c7-4-8}
If $x^{0}$ is a point for which
\[
\limsup_{R\to\infty}\left|\sum_{|m|<R}\left(1-\frac{|m|^{2}}{R^{2}}\right)^{(n-1)/2}e^{2\pi\mathrm{i}m\cdot x^{0}}\right|<\infty
\]

then
\begin{equation}\label{eq:c7-4-9}
\sup_{\substack{\alpha>(n-1)/2R>0}}\sup_{R>0}\left|\sum_{\lfloor m\rfloor<R}\left(1-\frac{\vert m\vert^{2}}{R^{2}}\right)^{\alpha}e^{2\pi i m\cdot x^{0}}\right|<\infty.
\end{equation}

\end{lemma}

This simple lemma indicates that it might be useful to reduce our consideration to the case $\alpha > (n - 1)/2$ where the Poisson summation formula is directly applicable.

To prove the lemma we begin by observing that the identity \footnote{This equality follows immediately, via a change of variables, from the well-known relationship between the beta and gamma functions: $\int_0^1 t^{x-1}(1 - t)^{y-1} dt = \Gamma(x) \Gamma(y) / \Gamma(x + y)$.}

\[
t^{\delta+\beta}=\frac{\Gamma(\delta+\beta+1)}{\Gamma(\delta+1)\Gamma(\beta)}\int\limits_{0}^{t}(t-s)^{\beta-1}s^{\delta}\,d s
\]

and the change of variables $s = r^{2} - |m|^{2}$ implies

\[
\begin{aligned}\int\limits_{\{m\}}^{R}(R^{2}-r^{2})^{\beta-1}\bigg(1&-\frac{|m|^{2}}{r^{2}}\bigg)^{(n-1)/2}r^{n} dr\\&=\int\limits_{\{m\}}^{R}(R^{2}-r^{2})^{\beta-1}(r^{2}-|m|^{2})^{(n-1)/2}r dr\\&=\frac{R^{2\beta+n-1}\Gamma[(n+1)/2]\Gamma(\beta)}{2\Gamma[(n+1)/2+\beta]}\left(1-\frac{|m|^{2}}{R^{2}}\right)^{\beta+(n-1)/2}\end{aligned}
\]

Putting

\[
\beta=\alpha-\frac{n-1}{2}>0,\qquad c_{\alpha}=\frac{2\Gamma(\alpha+1)}{\Gamma[(n+1)/2]\Gamma(\beta)}
\]

and interchanging the order of summation and integration we therefore obtain

\[
\begin{aligned}{K_{R}^{\alpha}(x)}&{{}=\sum_{|m|<R}\bigg(1-\frac{|m|^{2}}{R^{2}}\bigg)^{\alpha}e^{2\pi i m\cdot x}}\\ {}&{{}=\sum_{|m|<R}\Bigg\{c_{\alpha}R^{-2\alpha}\intop_{|m|}^{R}(R^{2}-r^{2})^{\beta-1}\bigg(1-\frac{|m|^{2}}{r^{2}}\bigg)^{(n-1)/2}r^{n}d r\Bigg\}e^{2\pi i m\cdot x}}\\ {}&{{}=c_{\alpha}R^{-2\alpha}\intop_{0}^{R}\bigg\{\sum_{|m|<r}\bigg(1-\frac{|m|^{2}}{r^{2}}\bigg)^{(n-1)/2}e^{2\pi i m\cdot x}\bigg\}(R^{2}-r^{2})^{\beta-1}r^{n}d r;}\\ \end{aligned}
\]

that is,

\begin{equation}\label{eq:c7-4-10}
K_{R}^{\alpha}(x)=c_{\alpha}R^{-2\alpha}\int\limits_{0}^{R}K_{r}(x)(R^{2}-r^{2})^{\beta-1}r^{n}d r.
\end{equation}

Now, if $\limsup_{R \to \infty} |K_R(x^0)| < \infty$ then $\sup_{R > 0} |K_R(x^0)| = A < \infty$. By \eqref{eq:c7-4-10} then

\[
|K_{R}^{\alpha}(x^{0})|\ \leq\ A\Bigg\{c_{\alpha}R^{-2\alpha}\int\limits_{0}^{R}(R^{2}\;-\;r^{2})^{\beta-1}r^{n}\ d r\Bigg\}
\]

This last expression within brackets, however, is identically equal to one (this can be easily seen from \eqref{eq:c7-4-10} by considering the constant term). Thus,

\[
\sup_{\substack{\alpha>(n-1)/2R>0}}\sup_{R>0}|K_{R}^{\alpha}(x^{0})|\leq\sup_{R>0}|K_{R}(x^{0})|
\]

and Lemma \ref{lem:c7-4-8} is proved.

We now describe the set $S$ of points for which (we claim) that \eqref{eq:c7-4-7} takes place. Let $S$ be the set of points $x$ where the denumerable collection of real numbers $\{|x - m|: m \in \Lambda\}$ is linearly independent over the rationals. In one dimension the set $S$ is obviously empty; however, we will show that the complement of $S$ has measure 0 when $n > 1$.

\begin{lemma}
\label{lem:c7-4-11}
If $x^0 \in S$, then
\[
\limsup_{R\to\infty}\left|\sum_{|m|<R}\left(1-\frac{|m|^{2}}{R^{2}}\right)^{(n-1)/2}e^{2\pi i m\cdot x^{0}}\right|=\infty.
\]

\end{lemma}

\textsc{Proof.}\ If $\alpha > (n-1)/2$, in view of the remarks made before \eqref{eq:c7-2-14}, we can apply the Poisson summation formula and obtain

\begin{equation}\label{eq:c7-4-12}
K_{R}^{\alpha}(x)=\pi^{-\alpha}\Gamma(\alpha+1)R^{(n/2)-\alpha}\sum_{m\in\Lambda}J_{(n/2)+\alpha}(2\pi R|x-m|)/|x-m|^{(n/2)+\alpha}.
\end{equation}

Using the asymptotic estimates for the Bessel functions in lemma 3.11 of Chapter IV (see also the accompanying footnote) we obtain for $x^0 \not\in \Lambda$

\begin{equation}\label{eq:c7-4-13}
K_{R}^{\alpha}(x^{0})=c_{\alpha}R^{(n-1)/2-\alpha}\Sigma\frac{\cos\left(2\pi R\gamma_{m}+\delta_{\alpha}\right)}{\gamma_{m}^{\alpha+(n+1)/2}}+E(R,\alpha)
\end{equation}

where $\gamma_m = |x^0 - m|$, and $\sup_{\alpha > (n-1)/2} \sup_R |E(R, \alpha)| < \infty$.

The identity \eqref{eq:c7-4-13} implies that

\begin{equation}\label{eq:c7-4-14}
\lim_{T\to\infty}\frac{1}{T}\int\limits_{1}^{T}K_{R}(x^{0})e^{2\pi i\lambda R}d R\text{exists}(=a(\lambda)),\text{if}\lambda\geq0,
\end{equation}

with $a(\lambda) = c\gamma_m^{-n}$, if $\lambda = \gamma_m$, and $a(\lambda) = 0$ otherwise. (To prove this integrate \eqref{eq:c7-4-13} with respect to $R$; then let $\alpha \to (n-1)/2$, and finally let $T \to \infty$). The following lemma shows that \eqref{eq:c7-4-14} is incompatible with $\limsup |K_R(x^0)| < \infty$, if $x^0 \in S$.

\begin{lemma}
\label{lem:c7-4-15}
Suppose $K$ is a bounded real function on $[1, \infty)$, and $\lim_{T \to \infty} \frac{1}{T} \int_{1}^{T} K(t) e^{2\pi i \lambda t} \, dt = a(\lambda)$ exists for every $\lambda \geq 0$. Let $\{\lambda_j\}$ be the set of $\lambda$'s for which $a(\lambda) \neq 0$, and assume this set is linearly independent over the
rationals. Then
\begin{equation}\label{eq:c7-4-16}
\Sigma|a(\lambda_{j})|<\infty.
\end{equation}

\end{lemma}

\textsc{Proof.}\ Let $\lambda_{1}, \ldots, \lambda_{n}$ be any finite sub-collection of $\{\lambda_{j}\}$. Write $a(\lambda_{j}) = |a(\lambda_{j})| e^{i\mu j}$, and $A(t) = \prod_{1}^{N} (1 + \cos (2\pi\lambda_{j}t - \mu_{j}))$. By the linear independence of $\{\lambda_{j}\}$, $\lim_{T \to \infty} \frac{1}{T} \int_{1}^{T} A(t) \, dt = 1$; moreover

\[
\lim_{T\to\infty}\frac{1}{T}\int\limits_{1}^{t}K(t)A(t)dt=\sum_{1}^{N}|a(\lambda_{j})|.
\]

From this it follows that $\sum_{1}^{n} |a(\lambda_j)| \leq \sup_{t \geq 1} |K(t)|$; hence \eqref{eq:c7-4-16} is proved, and so is lemma 4.15. Since

\[
\sum_{m}\gamma_{m}^{-n}\approx\sum_{m\neq0}|m|^{-n}=\infty,
\]

we see that $\limsup_{R \to \infty} |K_R(x^0)| = \infty$, if $x^0 \in S$. Therefore, Lemma \ref{lem:c7-4-11} is established and to complete the proof of our main lemma (Lemma \ref{lem:c7-4-6}) we need only show that the following result is true:

\begin{lemma}
\label{lem:c7-4-17}
If n > 1, the complement of S has measure 0.
\end{lemma}

\textsc{Proof.}\ Suppose $a_{m_1}, a_{m_2}, \ldots, a_{m_k}$ are nonzero rational numbers (associated with the lattice points $m_1, \ldots, m_k$). Let

\[
\Phi(x)=\sum_{j=1}^{k}a_{m_{j}}|x-m_{j}|.
\]

Then $\Phi(x) \not\equiv 0$ since $\Phi$ has singularities at $m_j, j = 1, 2, \ldots, k$. Moreover, $\Phi$ is real-analytic on the connected set $E_n - \Lambda$. But a nonzero real-analytic function can vanish only on a set of Lebesgue measure 0. Thus, the relation

\begin{equation}\label{eq:c7-4-18}
\Phi(x)=\sum_{j=1}^{k}a_{m_{j}}\vert x-m_{j}\vert=0
\end{equation}

can hold only in a set of measure 0. Since there are only a countable number of such point sets (each associated with a choice of nonzero rationals $a_{m_1}, \ldots, a_{m_k}$), their union also has measure 0. But this union is clearly $E_n - S.\footnote{The reader should note that this argument breaks down for $n = 1$ because $\Phi$ consists of distinct ``pieces'' of linear functions in this case; or to put it another way, $E_n - \Lambda$ is not connected, when $n = 1$.}$ This proves Lemma \ref{lem:c7-4-17}.

Let us see what we have accomplished. Let $\mu_{0}$ be the Dirac measure. Then

\[
d\mu_{0}\sim\sum_{m\in\Lambda}e^{2\pi i m\cdot x}.
\]

Lemma \ref{lem:c7-4-6}, therefore, shows that the Fourier--Stieltjes series of $d\mu_0$ is not summable by the Riesz means at the critical index for almost every $x \in Q_n$. To complete the proof of Theorem \ref{thm:c7-4-2} we shall replace $\mu_0$ by ``sharply peaked'' $L^1$ functions.

Let $\psi$ be a nonnegative $C^\infty$ function in $E_n$ with support in the unit ball $|x| \leq 1$ which satisfies $\int_{E_n} \psi(x) \, dx = 1$. Let $\Phi(y) = \int_{E_n} \psi(x) e^{-2\pi i x \cdot y} \, dx$ and $\varphi_\varepsilon(x) = \varepsilon^{-n} \sum_{m \in \Lambda} \psi((x - m)/\varepsilon)$. Then, by the Poisson summation formula,

\begin{equation}\label{eq:c7-4-19}
\varphi_{\varepsilon}(x)\sim\sum_{m\in\Delta}\Phi(\varepsilon m)e^{2\pi i m\cdot x}.
\end{equation}

We shall show that we can find two null sequences, $\{\varepsilon_k\}$ and $\{\delta_k\}$, of positive numbers so that the function $f$ defined by

\begin{equation}\label{eq:c7-4-20}
f(x)=\sum_{k=1}^{\infty}2^{-k}(\varphi_{\varepsilon_{k}}(x)-\varphi_{\delta_{k}}(x))
\end{equation}

belongs to $L^{1}(T_{n})$ and satisfies the conclusions of Theorem \ref{thm:c7-4-2}.

Let $S_{R}$ be the operator defined by

\[
\begin{aligned}(S_{R}f)(x)&=\sum_{|m|<R}\left(1\;-\;\frac{|m|^{2}}{R^{2}}\right)^{(n-1)/2}a_{m}e^{2\pi i m\cdot x}\\&\qquad\text{when}\;f\sim\sum_{m\in\Lambda}a_{m}e^{2\pi i m\cdot x}.\end{aligned}
\]

We first observe that

\begin{equation}\label{eq:c7-4-21}
\sup_{0<R}|(S_{R}\varphi_{\varepsilon})(x)|\leq A\varepsilon^{-n}.
\end{equation}

In fact,

\[
\sup_{0<R}|(S_{R}\varphi_{\varepsilon})(x)|\leq\sum_{m\in\Lambda}|\Phi(\varepsilon m)|=\sum_{|m|<1/\varepsilon}|\Phi(\varepsilon m)|+\sum_{|m|\geq1/\varepsilon}|\Phi(\varepsilon m)|.
\]

Since the functions with values $|\Phi(x)|$ and $|x|^{n+1}|\Phi(x)|$ are bounded, we have

\begin{equation}
\tag{i}
\sum_{\substack{|m|<1/\varepsilon}}\left|\Phi(\varepsilon m)\right|\leq\|\Phi\|_{\infty}\sum_{\substack{|m|<1/\varepsilon}}1=\|\Phi\|_{\infty}N_{\varepsilon},
\end{equation}

where $N_\varepsilon$ ( $\leq$ constant times $\varepsilon^{-n}$) is the number of lattice points satisfying $|m| < 1/\varepsilon$, and

\begin{equation}
\tag{ii}
\begin{aligned}\sum_{|m|\geq1/\varepsilon}\left|\Phi(\varepsilon m)\right|&\leq\left\{\sup_{x\in E_{n}}\left|x\right|^{n+1}\left|\Phi(x)\right|\right\}\sum_{|m|\geq1/\varepsilon}\left|\varepsilon m\right|^{-(n+1)}\\&=\left\{\sup_{x\in E_{n}}\left|x\right|^{n+1}\left|\Phi(x)\right|\right\}\varepsilon^{-n-1}\sum_{|m|\geq1/\varepsilon}\left|m\right|^{-(n+1)}\leq B\varepsilon^{-n}.\end{aligned}
\end{equation}

These two inequalities imply \eqref{eq:c7-4-21}.

We shall also construct subsets $\mathcal{E}_k \subset Q_n$ with measures $|\mathcal{E}_k| \geq 1 - 1/k$ and an increasing sequence of positive numbers $\{R_k\}$ such that

\begin{equation}\label{eq:c7-4-22}
\sup_{R\leq R_{k}}\left|(S_{R}f)(x)\right|\geq k,\qquad\text{if }x\in\mathcal{E}_{k}.
\end{equation}

Once this is done, it is easy to see that we have established Theorem \ref{thm:c7-4-2}. It is clear that the first conclusion of this theorem follows immediately from \eqref{eq:c7-4-22}. The second is a consequence of the fact that, in $Q_n$, $\varphi_\varepsilon$ has its support inside the solid sphere of radius $\varepsilon$ about the origin; therefore, by subtracting from $f$ partial sums of the series \eqref{eq:c7-4-20}, we obtain functions satisfying both the first and second parts of \eqref{thm:c7-4-2}.

Suppose now that $\varepsilon_j$, $\delta_j$, $R_j$, and $\mathcal{O}_j$ have been determined when $1 \leq j \leq k - 1$; we shall describe how to choose $\varepsilon_k$, $\delta_k$, $R_k$, and $\mathcal{O}_k$. We always select $\varepsilon_k \leq \delta_k$ and choose $\delta_k$ so small that

\begin{equation}\label{eq:c7-4-23}
\sup_{R\leq R_{k-1}}\left|S_{R}(\varphi_{\varepsilon_{k}}-\varphi_{\delta_{k}})\right|\leq1.
\end{equation}

If this can be done for each k, we have, of course,

\begin{equation}
\tag{4.23$'$}\label{eq:c7-4-23p}
\sup_{R\leq R_{k}}|S_{R}(\varphi_{\varepsilon_{k^{\prime}}}-\varphi_{\delta_{k^{\prime}}})|\leq1\qquad\mathrm{i f k^{\prime}>k.}
\end{equation}

Since

\[
\begin{aligned}{\sup_{R\leq R_{k-1}}\left|S_{R}(\varphi_{\varepsilon_{k}}-\varphi_{\delta_{k}})\right|}&{{}\leq\sum_{\substack{|m|<R_{k-1}}}|\Phi(\varepsilon_{k}m)-\Phi(\delta_{k}m)|}\\ {}&{{}\leq A(\delta_{k}-\varepsilon_{k})\sum_{\substack{|m|<R_{k-1}}}|m|\leq A^{\prime}\delta_{k}(R_{k-1})^{n+1},}\\ \end{aligned}
\]

we see that \eqref{eq:c7-4-23} can be realized if $\delta_{k}$ is small enough.

Now, let $A_{k}$ be a positive number such that

\begin{equation}\label{eq:c7-4-24}
A_{k}\geq\sup_{0<R<\infty}\Big\{\Big|S_{R}\Big(\sum_{j<k}2^{-j}(\varphi_{\varepsilon_{f}}-\varphi_{\delta_{f}})(x)-2^{-k}\varphi_{\delta_{k}}(x)\Big)\Big|\Big\}.
\end{equation}

Because of \eqref{eq:c7-4-21} such a finite constant $A_{k}$ can be found.

So far $A_k$ and $\delta_k$ have been fixed and $\epsilon_k$ is restricted only by the condition $\epsilon_k \leq \delta_k$. We now impose on $\epsilon_k$ a condition which, in essence, implies that $\varphi_{\epsilon_k}$ is sufficiently close to $d\mu_0$. That is, we know that if $R_k$ is large enough then, in view of the almost everywhere divergence in Lemma \ref{lem:c7-4-6},

\begin{equation}\label{eq:c7-4-25}
\sup_{R\leq R_{k}}2^{-k}|(S_{R} d\mu_{0})(x)|>A_{k}+k+2
\end{equation}

on a set $\mathcal{E}_k \subset Q_n$ having measure $|\mathcal{E}_k| \geq 1 - 1/k$. The Fourier series of $\varphi_\varepsilon$, however, converges termwise to that of $d\mu_0$ as $\varepsilon$ tends to 0. Therefore, by choosing $\varepsilon_k$ sufficiently small we have

\begin{equation}\label{eq:c7-4-26}
\sup_{R\leq R_{k}}2^{-k}|(S_{R}\varphi_{\varepsilon_{k}})(x)|>A_{k}+k+1\qquad\text{for}x\in\mathcal{E}_{k}.
\end{equation}

Thus, \eqref{eq:c7-4-25} and \eqref{eq:c7-4-26} determine our choice of $\varepsilon_{k}$, $R_{k}$, and $\mathcal{E}_{k}$.

We now turn to the proof of \eqref{eq:c7-4-22}. Write

\[
f=\Big\{\sum_{j<k}2^{-j}(\varphi_{\varepsilon_{1}}-\varphi_{\delta_{j}})-2^{-k}\varphi_{\delta_{k}}\Big\}+2^{-k}\varphi_{\varepsilon_{k}}+\Big\{\sum_{j>k}2^{-j}(\varphi_{\varepsilon_{j}}-\varphi_{\delta_{j}})\Big\}.
\]

Consider $\sup_{R \subseteq R_k} |(S_R f)(x)|$ for $x \in \mathcal{O}_k$. In view of \eqref{eq:c7-4-24} the contribution from the terms within the first bracket is at most $A_k$. By \eqref{eq:c7-4-26} the middle term contributes at least $A_k + k + 1$. Finally, the last term contributes at most 1 (by \eqref{eq:c7-4-23}). This proves \eqref{eq:c7-4-22} and Theorem \ref{thm:c7-4-2} is completely proved.

We shall now prove Theorem \ref{thm:c7-4-3}. We claim first that, if for some $x^{0}$

\[
\sup_{0<R<\infty}\bigg|\sum_{\substack{0<|m|<R}}\left|m\right|^{-(n/2)+1/2}e^{2\pi i m\cdot x^{0}}\bigg|<\infty,
\]

then

\begin{equation}\label{eq:c7-4-27}
\sup_{0<R<\infty}R^{-(n/2)+1/2}\left|\sum_{\substack{0<|m|<R}}e^{2\pi i m\cdot x^{0}}\right|<\infty.
\end{equation}

In fact, if we put $\sigma_R = \sum_{0 < |m| < R} |m|^{-(n/2)+1/2}e^{2\pi im \cdot x^0}$, then

\[
\begin{aligned}&\sum_{\substack{0<|m|<R}}e^{2\pi i m\cdot x^{0}}=\int\limits_{1}^{R}t^{(n/2)-1/2}d\sigma_{t}\\&=R^{(n/2)-1/2}\sigma_{R^{-}(n/2-1/2)}\int\limits_{1}^{R}\sigma_{t}t^{(n/2)-3/2}d t.\end{aligned}
\]

From this it is clear that the assumption $\sup_{0<R<\infty}|\sigma_R|<\infty$ implies \eqref{eq:c7-4-27}.

We next show that

\[
\sup_{0<R<\infty}R^{-(n/2)+1/2}\left|\sum_{0<|m|<R}e^{2\pi i m\cdot x^{0}}\right|=\infty
\]

for each $x^{0}$ satisfying

\[
\sup_{0<R<\infty}\left|\sum_{|m|<R}\left(1-\frac{|m|^{2}}{R^{2}}\right)^{(n-1)/2}e^{2\pi im\cdot x^{0}}\right|=\infty
\]

(by Lemma \ref{lem:c7-4-6}, almost all $x^0 \in Q_n$ satisfy this last relation). That is, we will show that

\[
\sup_{0<R<\infty}R^{-(n/2)+1/2}\big|K_{R}^{0}(x^{0})-1\big|=\infty
\]

for each $x^{0}$ satisfying

\[
\sup_{0<R<\infty}\;\big|K_{R}(x^{0})\big|=\sup_{0<R<\infty}\;\big|K_{R}^{(n-1)/2}(x^{0})\big|=\infty.
\]

It is easier to do this for $n$ odd and, in fact, the simplest case is the one occurring when $n = 3$. Let us consider this case; thus, the critical index is $(n - 1)/2 = 1$. For $x^{0}$ fixed write

\[
F_{\alpha}(t)=t^{\alpha}K_{\sqrt{t}}^{\alpha}(x^{0})=\sum_{|m|^{2}<t}(t-|m|^{2})^{\alpha}e^{2\pi i m\cdot x^{0}}.
\]

Obviously $(d/dt)F_\alpha(t) = \alpha F_{\alpha-1}(t)$ for $\alpha \geq 1$. We claim that

\begin{equation}\label{eq:c7-4-28}
|F_{2}(t)|\leq A t^{3/2}\qquad\text{for }1\leq t<\infty.
\end{equation}

In fact, when $\alpha > (n-1)/2$ the Poisson summation formula can be applied to $K_R^\alpha(x^0)$ in order to obtain the expression \eqref{eq:c7-4-12}. This, together with the estimate $|J_{(n/2)+\alpha}(2\pi R)| \leq AR^{-(1/2)}, R \geq 1$, (see Lemma \ref{lem:c4-3-11} in Chapter IV) shows that

\[
\big|K_{R}^{\alpha}\big(x^{0}\big)\big|\leq A R^{(n-1)/2-\alpha}\qquad\text{for }1\leq R<\infty
\]

which, by the definition of $F_{2}$, is equivalent to \eqref{eq:c7-4-28}.

If $|R^{-(n/2)+1/2}K_{R}^{0}(x^{0})| \leq A$ then (since $n = 3$)

\begin{equation}\label{eq:c7-4-29}
\frac{1}{2}\left|\frac{d^{2}F_{2}(t)}{d t^{2}}\right|=|F_{0}(t)|\leq A t^{1/2},\qquad1\leq t<\infty.
\end{equation}

Using the Taylor development of $F_{2}$ up to the second derivative, we see that

\[
|F_{2}(t+h)-F_{2}(t)-h F_{2}^{\prime}(t)|\leq\frac{h^{2}}{2}\sup_{t\leq t^{\prime}\leq t+h}|F_{2}^{\prime \prime}(t^{\prime})|.
\]

If we put $h = t^{1/2}$, \eqref{eq:c7-4-28} and \eqref{eq:c7-4-29} imply $|F_{2}^{\prime}(t)| \leq Bt$, $1 \leq t < \infty$. Thus, $|F_{1}(t)| \leq Bt/2$, and, therefore, $\sup_{1 \leq R < \infty} |K_{R}(x^{0})| \leq B$. This shows that

\[
\sup_{R>0}R^{-(n/2)+1/2}\big|K^{0}_{R}\big(x^{0}\big)-1\big|=\infty,
\]

whenever $\sup_{R>0}|K_{R}(x^{0})|=\infty$. This completes the proof of Thorem 4.3 in case $n=3$. The other cases $n>1$ are similar but somewhat more complicated. The needed machinery is presented in \eqref{frp:c7-6-10} below.

\section{Summability Below the Critical Index}

Suppose $\sum_{m \in \Lambda} a_m e^{2\pi im \cdot x}$ is the Fourier series of a function $f$ and

\[
S_{R}^{\alpha}(f)(x)=\sum_{|m|<R}\left(1-\frac{|m|^{2}}{R^{2}}\right)^{\alpha}a_{m}e^{2\pi i m\cdot x}
\]

are the corresponding Riesz means of order $\alpha$. We shall denote by $S_{*}^{\alpha}(f)(x)$ the associated ``maximal function''; that is,

\[
S_{*}^{\alpha}(f)(x)=\sup_{0<R<\infty}\left|S_{R}^{\alpha}(f)(x)\right|.
\]

In this section we shall see that if $f \in L^p(T_n)$, $1 < p < \infty$, then there are positive results for summability below the critical index. The theorem that is proved does not, in view of the counterexamples of Section 4, completely dispose of the problems of summability of multiple Fourier series.

Nevertheless, it does represent the furthest extent of our present knowledge of this subject. \footnote{See, however, the recent result of Fefferman [1].}

\begin{theorem}
\label{thm:c7-5-1}
Suppose $1 < p < \infty$, $n > 1$ and $\alpha > (n - 1)|^{\frac{1}{2}} - 1/p|$. Then
(a) $\|S_*(f)\|_p \leq A_{p,\alpha}\|f\|_p$;
(b) $\lim_{n \to \infty} S_n^\alpha(f)(x) = f(x)$ for almost every $x$;
(c) $\|S_R^\alpha(f) - f\|_p \to 0$ as $R \to \infty$.
\end{theorem}

This theorem will be deduced from two lemmas; the first will be a reformulation of the results obtained in the second section when $\alpha$ was greater than the critical index, the second will be an $L^{2}$ result for $\alpha$ near zero and, finally, a combination of these two estimates obtained by using the complex convexity of the norms of linear operators (Theorem \ref{thm:c5-4-1} in Chapter V). The first lemma is as follows:

\begin{lemma}
\label{lem:c7-5-2}
When $\alpha > (n-1)/2$ conclusion (a) of Theorem \ref{thm:c7-5-1} is true.
\end{lemma}

We remind the reader that when $\alpha > (n-1)/2$, conclusions (b) and (c) of Theorem \ref{thm:c7-5-1} are contained in Theorem \ref{thm:c7-2-11}, and in particular its Corollary \ref{cor:c7-2-15} above.

\textsc{Proof.}\ Let $f \in L^p(T_n)$. We also use the symbol $f$ to denote the periodic extension to $E_n$. By Corollary \ref{cor:c7-2-15} (see also the comments preceding it) we know that

\[
S_{R}^{\alpha}(f)(x)=\int\limits_{Q_{n}}\sum_{m\in\Lambda}\varphi_{\varepsilon}(y+m)f(x-y)d y=\int\limits_{E_{n}}\varphi_{\varepsilon}(y)f(x-y)d y
\]

for an appropriate $\varphi$, where $\varphi_{\varepsilon}(x) = \varepsilon^{-n}\varphi(x/\varepsilon)$ and $\varepsilon = 1/R$.

Observe that in proving estimate (a) for $S_{*}^{\alpha}(f) = \sup_{0 < R < \infty} |S_{R}^{\alpha}(f)|$, it suffices to restrict our attention to $R \geq 1$, (i.e., $\varepsilon \leq 1$) since when $R < 1$ then $S_{R}^{\alpha}(f) = a_0$, and $|a_0| \leq \|f\|_1 \leq \|f\|_p$. Now let $\tilde{f}(x) = 0$ if $|x| > 1$, and $\tilde{f}(x) = f(x)$, if $|x| \leq 1$. Then, clearly,

\[
\biggl(\intop_{E_{n}}|\tilde{f}(x)|^{p}d x\biggr)^{1/p}\leq A\biggl(\intop_{Q_{n}}|f(x)|^{p}d x\biggr)^{1/p}.
\]

Also

\[
(f*\varphi_{\varepsilon})(x)=\int\limits_{|y|\leq1}\varphi_{\varepsilon}(x-y)f(y)d y+\int\limits_{|y|>1}\varphi_{\varepsilon}(x-y)f(y)d y.
\]

We have already observed (see the paragraph preceding \eqref{eq:c7-2-14}) that, in this case,

\[
\begin{array}{r}{\varphi(x)=\pi^{-\alpha}\Gamma(\alpha+1)|x|^{-(\alpha+n/2)}J_{\alpha+n/2}(2\pi|x|).}\end{array}
\]

Thus, by \eqref{eq:c4-3-12} of Chapter IV, $|\varphi(x)| \leq A(1 + |x|)^{-n-\delta}$ with $\delta = \alpha - (n-1)/2$. Consequently, as was shown in the proof of Theorem \eqref{thm:c7-2-11},

\[
\bigg|\int_{|y|>1}\varphi_{\varepsilon}(x-y)f(y)\;d y\bigg|\;\leq\;\mathcal A^{\prime}\varepsilon^{\delta}\bigg(\sum_{m\neq0}|m|^{-n-\delta}\bigg)\int_{Q_{n}}|f(x)|\;d x.
\]

But

\[
\int\limits_{|y|\leq1}\varphi_{\varepsilon}(x-y)f(y)d y=\int\varphi_{\varepsilon}(x-y)\tilde{f}(y)d y
\]

and, therefore, by \eqref{eq:c2-3-9} of Chapter II,

\[
\Vert S_{*}^{\alpha}(f)\Vert_{p}\;\leq\;A_{p,\alpha}\Vert f\Vert_{p}
\]

if p > 1 and $\alpha > (n - 1)/2$.

Let us make some general remarks concerning the Riesz means $S_R^a$. When we introduced them we made the implicit assumption that the index $\alpha$ is nonnegative. Nevertheless, nothing in the definition we gave impedes our considering these means when $\alpha$ is negative or even complex. In fact, the behavior of $S_R^a$ as an analytic function of $\alpha$ will be one of the main tools we shall use. In the following $L^2$ result the means of negative order are already involved.

\begin{lemma}
\label{lem:c7-5-3}
Let
\[
M^{\alpha}(f)(x)=\sup_{0<R<\infty}\left(\frac{1}{R}\int_{0}^{R}|S_{t}^{\alpha}(f)(x)|^{2}\,d t\right)^{1/2}.
\]

\end{lemma}

Then

\begin{equation}\label{eq:c7-5-4}
\|M^{\alpha}(f)\|_{2}\leq A_{\alpha}\|f\|_{2}
\end{equation}

if $f \in L^2(T_n)$ and $\alpha > -\tfrac{1}{2}$. \footnote{Notice that the quantities $M^{\alpha}(f)$ represent averages of the type of expressions we want eventually to dominate.}

\textsc{Proof.}\ We introduce the auxiliary function $G^{\alpha}(f)$ defined by

\[
G^{\alpha}(f)(x)=\left(\int\limits_{0}^{\infty}\left|S_{R}^{\alpha+1}(f)(x)-S_{R}^{\alpha}(f)(x)\right|^{2}\frac{d R}{R}\right)^{1/2}.
\]

We claim that

\begin{equation}\label{eq:c7-5-5}
\left\|{G}^{\alpha}(f)\right\|_{2}\leq A_{\alpha}\|f\|_{2}
\end{equation}

if $f \in L^2(T_n)$ and $\alpha > -\frac{1}{2}$. In fact, by Parseval's relation and Fubini's theorem

\[
\begin{aligned}\|G^{\alpha}(f)\|_{2}^{2}&=\int\limits_{0}^{\infty}\left(\sum_{|m|<R}\left|\left(1-\frac{|m|^{2}}{R^{2}}\right)^{(\alpha+1)}-\left(1-\frac{|m|^{2}}{R^{2}}\right)^{\alpha}\right|^{2}|a_{m}|^{2}\right)\frac{dR}{R}\\&=\sum_{m\neq0}|a_{m}|^{2}\int\limits_{|m|}^{\infty}\frac{|m|^{4}}{R^{4}}\left[1-\frac{|m|^{2}}{R^{2}}\right]^{2\alpha}\frac{dR}{R}=C_{\alpha}\sum_{m\neq0}|a_{m}|^{2}.\end{aligned}
\]

The third equality is a consequence of the fact that

\[
\begin{aligned}\int\limits_{|m|}^{\infty}\frac{|m|^{4}}{R^{4}}\left[1-\frac{|m|^{2}}{R^{2}}\right]^{2\alpha}\frac{dR}{R}&=\int\limits_{1}^{\infty}s^{-5}(1-s^{-2})^{2\alpha}ds=C_{\alpha}\\&=[2(2\alpha+1)(2\alpha+2)]^{-1}\end{aligned}
\]

and the last integral converges if $\alpha > -\frac{1}{2}$. Therefore,

\[
\|G^{\alpha}(f)\|_{2}^{2}=C_{\alpha}\{\|f\|_{2}^{2}-|a_{0}|^{2}\},
\]

which is, of course, a more precise form of \eqref{eq:c7-5-5}. From \eqref{eq:c7-5-5} the lemma follows easily. To begin with, notice that by the definition of $G^{\alpha}(f)$

\[
\sup_{0<R<\infty}\frac{1}{R}\int\limits_{0}^{R}\left|S_{t}^{\alpha+1}(f)(x)-S_{t}^{\alpha}(f)(x)\right|^{2}d t\leq[G^{\alpha}(f)(x)]^{2},
\]

and, thus

\begin{equation}\label{eq:c7-5-6}
M^{\alpha}(f)(x)\leq M^{\alpha+1}(f)(x)+G^{\alpha}(f)(x).
\end{equation}

Consequently, applying \eqref{eq:c7-5-6} repeatedly,

\begin{equation}\label{eq:c7-5-7}
M^{\alpha}(f)\leq M^{\alpha+k}(f)+G^{\alpha}(f)+G^{\alpha+1}(f)+\cdot\cdot\cdot+G^{\alpha+k-1}(f).
\end{equation}

Now, if we choose $k > n/2$ then $\alpha + k > (n-1)/2$. Since $M^{\alpha+k}(f)$ is clearly less than or equal to $S_{*}^{\alpha+k}(f)$, we can apply Lemma \ref{lem:c7-5-2} (with $p = 2$) and obtain $\|M^{\alpha+k}(f)\|_{2} \leq A\|f\|_{2}$. This inequality, together with \eqref{eq:c7-5-5} (applied to $G^{\alpha}(f)$, $G^{\alpha+1}(f)$,..., $G^{\alpha+k-1}(f)$) gives us \eqref{eq:c7-5-4} and the lemma is proved.

We can use Lemma \ref{lem:c7-5-3} once we have obtained the formula which expresses the Riesz means of a given order as averages of the Riesz means of lower orders. The required calculation was already carried out in the proof of Lemma \ref{lem:c7-4-8}. In fact, the identity

\[
\Bigl(1-\frac{|m|^{2}}{R^{2}}\Bigr)^{\beta+\delta}=C_{\beta,\delta}R^{-2\beta-2\delta}\intop_{1\leq i\leq n}^{R}(R^{2}-t^{2})^{\beta-1}t^{2\delta+1}\Bigl(1-\frac{|m|^{2}}{t^{2}}\Bigr)^{\delta}d t,
\]

with $C_{\beta,\delta} = 2\Gamma(\delta + \beta + 1)/\Gamma(\delta + 1)\Gamma(\beta)$, gives us immediately the desired expression; namely,

\begin{equation}\label{eq:c7-5-8}
S_{R}^{\beta+\delta}=C_{\beta,\delta}R^{-2\beta-2\delta}\int_{0}^{R}(R^{2}-t^{2})^{\beta-1}t^{2\delta+1}S_{t}^{\delta}d t.
\end{equation}

It then follows that

\[
\begin{aligned}{|S_{R}^{\beta+\delta}(f)(x)|\leq}&{{}C_{\beta,\delta}R^{-2\beta-2\delta}\Bigg(\int_{0}^{R}|(R^{2}-t^{2})^{\beta-1}t^{2\delta+1}|^{2}d t\Bigg)^{1/2}}\\ {}&{{}\times R^{1/2}\cdot R^{-1/2}\left(\int_{0}^{R}|S_{t}^{\delta}(f)(x)|^{2}d t\right)^{1/2};}\\ \end{aligned}
\]

therefore, upon taking the supremum over all R > 0, we obtain

\begin{equation}\label{eq:c7-5-9}
S_{*}^{\beta+\delta}(f)(x)\leq C_{\beta,\delta}^{\prime}M^{\delta}(f)(x)\qquad\text{if }\beta>\frac{1}{2}.\footnote{Notice that, if $\beta > \frac{1}{2}$, then 
\[
R^{1-4\beta-4\delta}\int\limits_{0}^{R}\vert(R^{2}-t^{2})^{\beta-1}t^{2\delta+1}\vert^{2}\;{d}t=\int\limits_{0}^{1}\vert(1-t^{2})^{\beta-1}t^{2\delta+1}\vert^{2}\;{d}t<\infty.
\]}
\end{equation}

Finally, given $\alpha > 0$, choose $\beta$ and $\delta$ so that $\beta + \delta = \alpha$, with $\delta > -\frac{1}{2}$ and $\beta > \frac{1}{2}$. Then, combining \eqref{eq:c7-5-9} and Lemma \ref{lem:c7-5-3}, we obtain the following basic inequality in the $L^{2}$ case:

\begin{lemma}
\label{lem:c7-5-10}
If $\alpha > 0$ then $\|S_*(f)\|_2 \leq A_\alpha\|f\|_2$ for all $f$ in $L^2(T_n)$.
\end{lemma}

Let us now pause momentarily in our efforts toward proving Theorem \ref{thm:c7-5-1} in order to see what has already been accomplished. The principal part of that theorem is the maximal inequality (part (a)); once it is established, the rest is routine. When $p$ is near 1 or $\infty$, part (a) is in effect already contained in Lemma \ref{lem:c7-5-2}. Moreover, Lemma \ref{lem:c7-5-10} gives us the desired inequality in the special case $p = 2$. The general case can be deduced from these special cases by applying the interpolation theorem that was established in the fourth section of Chapter V. Although the ideas we use in doing this are simple, the details required are somewhat complicated. What is needed first is a strengthening of the two special cases just mentioned which allows orders $\alpha$ that are complex. This strengthening is a consequence of a very simple general principle which asserts that if we have an estimate for $S_{*}^{\alpha}(f)$ with $\alpha \geq 0$, then there exists a corresponding estimate for $S_{*}^{\alpha}(f)$ with $\alpha^{\prime}$ complex and satisfying $\text{Re}(\alpha^{\prime}) > \alpha$. That this is so follows from formula \eqref{eq:c7-5-8} with $\delta = \alpha, \beta + \delta = \alpha^{\prime}$ (and, thus

Re $\left(\beta\right)>0$. In fact, we have

\[
\begin{aligned}S_{*}^{\alpha^{\prime}}(f)&\leq S_{*}^{\alpha}(f)C_{\beta,\delta}\int\limits_{0}^{1}(1-s^{2})^{\text{Re}(\beta)-1}s^{2\alpha+1}ds\\&=\frac{|\Gamma(\alpha^{\prime}+1)|\Gamma(\text{Re}(\beta))}{\Gamma(\text{Re}(\alpha^{\prime})+1)|\Gamma(\beta)|}S_{*}^{\alpha}(f).\end{aligned}
\]

Since $\Gamma(\alpha' + 1) \leq |\Gamma(Re(\alpha') + 1)|$ this gives us

\begin{equation}\label{eq:c7-5-11}
S_{*}^{\alpha^{\prime}}(f)\;\leq\;D_{\beta}S_{*}^{\alpha}(f)
\end{equation}

with $D_{\beta} = \Gamma(\text{Re}(\beta))/|\Gamma(\beta)|$.

Since the operators $f \to S_*(f)$ are not linear, and we wish to apply an interpolation theorem for linear operators, we need to introduce another technical device. In order to do this we let $\mathcal{C}$ denote the class of nonnegative measurable functions on $T_n$ having only a finite number of distinct values. Suppose $R \in \mathcal{C}$, then

\begin{equation}\label{eq:c7-5-12}
|S_{R(x)}^{\alpha}(f)(x)|\leq S_{*}^{\alpha}(f)(x).
\end{equation}

We assert, however, that the following converse of this inequality holds:

\begin{equation}\label{eq:c7-5-13}
\sup_{R\in\mathcal{C}}\|S_{R(x)}^{\alpha}(f)(x)\|_{p}=\|S_{*}^{\alpha}(f)(x)\|_{p}.
\end{equation}

This follows from the obvious fact that we can find a sequence $R_1(x)$, $R_2(x)$,..., $R_j(x)$,... in C such that

\[
\lim_{j\to\infty}\left|S_{R_{j}(x)}^{\alpha}(f)(x)\right|=\sup_{0<R<\infty}\left|S_{R}^{\alpha}(f)(x)\right|=S_{*}^{\alpha}(f)(x)\qquad\text{for all }x\in T_{n}.
\]

We now fix an element $R \in \mathcal{C}$ and consider the linear operators

\begin{equation}\label{eq:c7-5-14}
f\to S_{R(x)}^{\alpha}(f)(x)
\end{equation}

and their analytic dependence on the parameter $\alpha$. Suppose $f$ is integrable and $f(x) \sim \sum_{m \in \Lambda} a_m e^{2\pi i m \cdot x}$. Since $R(x)$ takes on only a finite number of values (in particular, it is bounded) it is easy to see that

\[
\begin{array}{r}{S_{R(x)}^{\alpha}(f)(x)=\displaystyle\sum_{|m|\le R(x)}\left(1\;-\;\frac{|m|^{2}}{R^{2}(x)}\right)^{\alpha}a_{m}e^{2\pi i m\cdot x}}\end{array}
\]

is analytic in $\alpha$ and is of admissible growth in any vertical strip of the complex $\alpha$ plane, in the sense of Theorem \ref{thm:c5-4-1} in Chapter V.

Let $\alpha = \mu + i\nu$, $\mu_0 > 0$ and $\mu_1 > (n-1)/2$. Then, by \eqref{eq:c7-5-12}, \eqref{eq:c7-5-11}, and \eqref{lem:c7-5-10} with $\alpha' = \mu_0 + i\nu$, $\alpha = \mu_0/2$, $\beta = (\mu_0/2) + i\nu$, we have

\begin{equation}\label{eq:c7-5-15}
\left\|S_{R(x)}^{\mu_{0}+i\nu}(f)(x)\right\|_{2}\leq A_{0}(\nu)\left\|f\right\|_{2}
\end{equation}

for all simple functions f, where

\[
A_{0}(\nu)\leq A_{0}\;\frac{\Gamma(\mu_{0}/2)}{|\Gamma((\mu_{0}/2)\;+\;i\nu)|}\leq A_{0}^{\prime}e^{\pi|\nu|}.\footnote{More precise estimates for $A_{0}(\nu)$ can be made by using the asymptotic formula $\left|\Gamma((\mu/2)+i\nu)\right| \sim \sqrt{2\pi}|\nu|^{(\mu-1)/2} e^{-\pi|\nu|/2}$, as $\nu \to \infty$, but they are not needed in order to apply the interpolation Theorem \ref{thm:c5-4-1} of Chapter V.}
\]

What is important to observe is that the bound for $A_0(\nu)$ does not depend on the particular choice of $R\in\mathcal{C}$. Similarly, by \eqref{eq:c7-5-12}, \eqref{eq:c7-5-11}, and \eqref{lem:c7-5-2} with $\alpha'=\mu_1+i\nu$, $\alpha=\frac{1}{2}[\mu_1+(n-1)/2]$, $\beta=\frac{1}{2}[\mu_1-(n-1)/2]+i\nu$, we have

\begin{equation}\label{eq:c7-5-16}
\left\|S_{R(x)}^{\mu_{1}+\mathrm{i}\nu}(f)(x)\right\|_{p_{1}}\leq A_{1,p_{1}}(\nu)\left\|f\right\|_{p_{1}}
\end{equation}

for all simple $f, 1 < p_1 < \infty, [\mu_1 > (n-1)/2]$, where $A_{1,p_1}(\nu) \leq A_{1,p_1}' e^{\pi |\nu|}$. Here again the constants are independent of $R(x)$.

We now invoke the interpolation Theorem \ref{thm:c5-4-1} of Chapter V. In keeping with the notation set down there we let $0 < t < 1$, $p_0 = 2$ and $p_1$ be the exponent arising in \eqref{eq:c7-5-16}. Then, if $\mu = \mu_0(1 - t) + \mu_1 t$ and $1/p = [(1 - t)/p_0] + [t/p_1] = (1 - t)/2 + t/p_1$, we obtain

\begin{equation}\label{eq:c7-5-17}
\left\|S_{R(x)}^{\mu}(f)(x)\right\|_{p}\leq A_{p}\left\|f\right\|_{p}.
\end{equation}

Again $A_{p}$ is independent of the $R(x)$ chosen from C. But \eqref{eq:c7-5-17} and \eqref{eq:c7-5-13} give us the inequality:

\begin{equation}\label{eq:c7-5-18}
\|S_{*}^{\mu}(f)\|_{p}\leq A_{p}\|f\|_{p}.
\end{equation}

We claim that the restrictions on $p$ and $\mu$ in \eqref{eq:c7-5-18} are precisely those of Theorem \ref{thm:c7-5-1}; that is, $1 < p < \infty$ and $\mu > (n-1)|1\frac{1}{2} - 1/p$. Recall that $\mu_0 > 0$, $\mu_1 > (n-1)/2$ and $1 < p_1 < \infty$ are otherwise arbitrary, while $\mu = \mu_0(1-t) + t\mu_1$ and $1/p = (1-t)/2 + t/p_1$. Assume first that $p \leq 2$. If we set $p_1 = 1$, $\mu_0 = 0$, $\mu_1 = (n-1)/2$, then an obvious calculation gives us $\mu = (n-1)[1/p-\frac{1}{2}]$. It is also clear that $\mu$ is a continuous function of $p_1$, $\mu_0$, and $\mu_1$ (we see this by expressing $t$ in terms of these parameters). Thus, by continuity, we can always realize any $\mu$ satisfying $\mu > (n-1)(1/p-\frac{1}{2})$ by choosing $p_1 > 1$, $\mu_0 > 0$ and $\mu_1 > (n-1)/2$. If $p \geq 2$ the argument is similar once we begin by setting $p_1 = \infty$. This concludes the proof of part (a) of Theorem \ref{thm:c7-5-1}.

Parts (b) and (c) follow from this by the general principles concerning convergence of families of operators developed in the second chapter (see, in particular, Theorem \ref{thm:c2-3-12}) when we observe that, if we choose $f$ from the class of trigonometric polynomials (which is dense in $L^{p}(T_{n})$ when $p < \infty$), then $\lim_{R \to \infty} S_{R}^{\alpha}(f)(x)$ converges uniformly to $f$, for each $\alpha \geq 0$.

\section{Further Results}

\frp{6}{1} An application of reasoning similar to that used in the proof of Theorem \ref{thm:c7-2-17} shows that each of the following are Fourier series of functions that are continuously differentiable and periodic on the fundamental cube except at the origin; moreover, they have the indicated behavior:

(a) $\sum_{m \neq 0} |m|^{-n} e^{2\pi im \cdot x}$ is asymptotic to $c \cdot \log (1/|x|)$ as $|x| \to 0$;

(b) If $P_k$ is a homogeneous harmonic polynomial of degree $k \geq 1$, then $\sum_{m \neq 0} |m|^{-n-k}P_k(m)e^{2\pi im \cdot x}$ is a bounded function;

(c) $\sum_{m \neq 0} |m|^{-\alpha-k} b(|m|) P_k(m) e^{2\pi im \cdot x}$ is asymptotic to $c|x|^{-(n-\alpha)-k} \times b(1/|x|) P_k(x)$ as $|x| \to 0$, when $P_k$ is a homogeneous harmonic polynomial of degree $k \geq 0$, $0 < \alpha < n$ and $b(x)$ is an appropriate function slowly varying near $\infty$; (examples of allowable $b$'s are functions whose values for $|x|$ large are $(\log|x|)^a$, $(\log(\log|x|))^c$, $\exp(\log|x|)^{1/2}$, etc., and also products and quotients of such functions). For the one-dimensional theory see Zygmund [1] and for the $n$-dimensional one see Wainger [1].

\frp{6}{2} Some examples of a different nature may also be obtained by the aid of the Poisson summation formula:

(a) $\sum_{m \neq 0} e^{ci|m|\log|m|}|m|^{-\varepsilon-n/2}e^{2\pi im\cdot x}$ is the Fourier series of a continuous periodic function, if $c \neq 0$ and $\varepsilon > 0$;

(b) If the dimension $n$ is even then $\sum_{|m|>1}|m|^{-n}(\log|m|)^{-1} \times e^{ci|m|\log(|m|)^a}e^{2\pi im \cdot x}$, with $c \neq 0$ and $0 < a < 2/n$, is the Fourier series of a function of class $C^{n/2}$, but is not absolutely convergent.

In the one-dimensional case example (a) is due to Hardy and Littlewood. For both (a) and (b) see Wainger [1], where other examples are also studied. Example (b), incidentally, shows that the simple Corollary \ref{cor:c7-1-9} cannot be improved very much.

\frp{6}{3} The following considerations give an interpretation of the heuristic formula \eqref{eq:c7-2-16}. Let $P_k(x)$ be a homogeneous harmonic polynomial of degree k. Fix x and consider the following two functions of $\alpha$:

\[
\sum_{m\neq0}\frac{P_{k}(m)}{|m|^{k+\alpha}}e^{2\pi i m\cdot x},\qquad\text{for Re }(\alpha)>n,
\]

and

\[
\gamma_{\alpha,k}^{-1}\sum_{m\neq0}\frac{P_{k}(x+m)}{|x+m|^{k+n-\alpha}},\qquad\text{for Re }(\alpha)<0,
\]

where $\gamma_{\alpha,k} = i^{-k}\pi^{n/2-\alpha}\Gamma((k + \alpha)/2)/\Gamma((n + k - \alpha)/2)$. These two functions of $\alpha$ are equal in the sense that they are analytic continuations of each other (see Hecke [1] and Bochner [6]).

\frp{6}{4} Let $S \subset E_n$ be convex, symmetric, open and bounded. Suppose $|S| > 2^n$. Then $S$ contains at least one lattice point different from the origin. This classical theorem of Minkowski can be proved by making use of the Poisson summation formula in the following way. Let $S_{1/2} = \{x: 2x \in S\}$ and let $\varphi$ denote the characteristic function of $S_{1/2}$. We set $f = \varphi * \varphi$; thus, $\hat{f} = |\hat{\varphi}|^2 \geq 0$. It can be shown that if the origin is the only lattice point in $S$ then $f(m) = 0$ when $m \neq 0$. Moreover, the Poisson summation formula $\sum f(m) = \sum \hat{f}(m)$ holds. If 0 is the only lattice point within $S$, therefore, $f(0) = \sum \hat{f}(m) \geq \hat{f}(0)$. But this contradicts the fact that $f(0) = \int \varphi(x) \, dx = |S_{1/2}|$ and $\hat{f}(0) = \int f(x) \, dx = |S_{1/2}|^2$. The argument sketched goes back to Siegel [1].

\frp{6}{5} Suppose $\sum_{m\in\Lambda} a_m e^{2\pi im \cdot x}$ is a given formal trigonometric series. The $n$ series

\[
-i\sum_{m\neq0}a_{m}\frac{m_{k}}{|m|}e^{2\pi i m\cdot x},\qquad k=1,\ldots,n,
\]

are the Riesz transforms of the given series. Let

\[
u_{0}(x,t)=\sum_{m\in\Lambda}a_{m}e^{-2\pi\vert m\vert t}e^{2\pi i m\cdot x}
\]

and

\[
u_{k}(x,t)=-i\sum_{m\neq0}a_{m}\frac{m_{k}}{|m|}e^{-2\pi|m|t}e^{2\pi i m\cdot x}\qquad\text{for }t>0
\]

(assuming that the series defining $u_k, k = 0, \ldots, n$, converges). We observe that the $(n + 1)$-tuple $F = (u_0, u_1, \ldots, u_n)$ satisfies the generalized Cauchy-Riemann equations \eqref{eq:c6-4-13} of Chapter VI. In analogy with the theory described there we say that $F \in H^p(T_n)$ if $\sup_{t > 0} \int_{T_n} |F(x, t)|^p \, dx < \infty$.

(a) $F \in H^p(T_n)$, $1 < p < \infty$, if and only if $\sum_{m \in \Lambda} a_m e^{2\pi i m \cdot x}$ is the Fourier series of a function in $L^p(T_n)$. This can be proved by using the results developed in Section 3.

(b) Suppose $(n - 1)/n < p < \infty$ and $F \in H^p(T_n)$, then

\[
\int\limits_{T_{n}}\sup_{t>0}|F(x,t)|^{p}dt<\infty
\]

and $\lim F(x, t)$ exists almost everywhere and in the metric of $L^{p}(T_{n})$.

(c) The special case $p = 1$ implies the following extension of the classical F. and M. Riesz theorem (see also \eqref{frp:c6-5-8} in Chapter VI). Suppose $\sum_{m \in \Lambda} a_m e^{2\pi i m \cdot x}$ and $\sum_{m \neq 0} a_m (m_k / |m|) e^{2\pi i m \cdot x}$, $k = 1, 2, \ldots, n$. are Fourier-Stieltjes series of finite measures, then these measures are absolutely continuous.

Both (b) and (c) can be established by arguments that are analogous to those presented in Chapter VI. See also Shapiro [1].

(d) A result that complements the last result (which has no analog when $n=1$) is as follows. Suppose that $\sum_{m\in\Lambda} a_m e^{2\pi i m\cdot x}$ and $\sum_{m\neq0} a_m(m_k/|m|) e^{2\pi i m\cdot x}$, $k=1,2,\ldots,n$, are Fourier-Stieltjes series of finite measures, then for each homogeneous (of degree $r$) harmonic polynomial $P_r(x)$, $r=0,1,2,\ldots$, the series $\sum_{m\neq0}(P_r(m)/|m|^r) a_m e^{2\pi i m\cdot x}$ is the Fourier series of a function in $L^1(T_n)$. This is a consequence of techniques presented in Stein [3], Chapter VII.

\frp{6}{6} Theorem \ref{thm:c7-2-11} and Lemma \ref{lem:c7-5-2} have the following more general formulation. Suppose $\varphi \in L^1(E_n)$ with $\int_{E_n} \varphi(x) \, dx = 0$. Let $\varphi_\varepsilon(x) = \varepsilon^{-n} \varphi(x/\varepsilon)$ and $K_\varepsilon(x) = \sum_{m \in \Lambda} \varphi_\varepsilon(x + m)$. Consider $(f * K_\varepsilon)(x) = \int_{T_n} f(x - y) K_\varepsilon(y) \, dy$. Then

(a) If $f \in L^p(T_n)$, $1 \leq p < \infty$, then $f \cdot K_\varepsilon \to f$ in the $L^p(T_n)$ norm as $\varepsilon \to 0$;

(b) Suppose, in addition, that $\psi(x) = \sup_{|x| \leq |x'|} |\varphi(x')|$ is integrable on $E_n$. Then $(f * K_\varepsilon)(x) \to f(x)$ at each point of the Lebesgue set of $f$. Moreover, if $\tilde{f}(x) = f(x)$ for $|x| \leq 1$ and $\tilde{f}(x) = 0$ if $|x| > 1$ then

\[
\sup_{\varepsilon>0}\left|(f* K_{\varepsilon})(x)\right|\leq A m_{\tilde{f}}(x).
\]

\frp{6}{7} Let $\lambda$ be a complex-valued function defined on $E_n$ and of class $C^n$ except at the origin. Suppose that there exists a constant $A$ such that $|\mathcal{D}^\alpha \lambda(x)| \leq A/|x|^{|\alpha|}$ whenever $0 \leq |\alpha| \leq n$. Then the sequence $\{\lambda(m)\}, m \in \Lambda, (\lambda(0) = 0)$ is a multiplier sequence of type $(L^p(T_n), L^p(T_n))$ for $1 < p < \infty$. This theorem goes back to Marcinkiewicz [1], where a somewhat stronger assertion is proved. Under the same conditions, $\lambda$ is also a multiplier in the nonperiodic case (i.e., $\lambda \in (L^p(E_n), L^p(E_n))$, for $1 < p < \infty$), but for historical reasons this type of result appeared only much later (see Mihlin [1] and Hörmander [1]). The nonperiodic version could also be deduced directly from the periodic version by making use of Theorem \ref{thm:c7-3-18}. See also Stein [3], Chapter IV.

\frp{6}{8} Theorem \ref{thm:c7-5-1} shows that $\|S_*(n-1)/2(f)\|_p \leq A_p \|f\|_p, 1 < p < \infty$, if $f \in L^p(T_n)$. This result for the critical index may be extended by considering the growth of $A_p$ as $p \to 1$ or $p \to \infty$. In fact, a refinement of the proof of Theorem \ref{thm:c7-5-1} shows that $A_p \leq A(p-1)^{-2}$ for $1 < p < 2$ and $A_p \leq Ap$ for $2 \leq p < \infty$.

(a) In view of the estimate for $A_p$ for $p$ near 1 we obtain the inequality $\|S_*(^{n-1)/2}(f)\|_1 \leq A \int_{T_n} |f|(\log^+ |f|)^2 dx + B$. Consequently, if $|f|(\log^+ |f|)^2$ is integrable, $\lim_{R \to \infty} S_R^{(n-1)/2}(f)(x) = f(x)$ for almost every $x$.

(b) From the estimate for $A_p$ for large $p$ it follows that if $f$ is bounded then there exists $a > 0$ such that $\int_{T_n} \exp \{a S_{*}^{(n-1)/2}(f)(x)\} \, dx < \infty$.

For part (a) see Stein [4]; the estimate for $A_{p}$ needed to establish (b) is implicitly contained there.

\frp{6}{9} We have observed that localization does not hold for arbitrary functions in $L^{1}(T_{n})$ at the critical index (see Theorem \ref{thm:c7-4-2} and the accompanying footnote). Localization, however, does hold under stronger assumptions. If we suppose $\int_{Q_{n}} |f| \log^{+} |f| \, dx < \infty$, for example, localization does hold. In particular, if at a given point $x_{0}$, $f$ satisfies the Dini condition $\int_{Q_{n}} |f(x_{0}-t) - f(x_{0})| \, |t|^{-n} \, dt < \infty$ then $S_{R}^{(n-1)/2}(f)(x_{0}) \to f(x_{0})$ as $R \to \infty$. This is a consequence of the following result:

\[
\sup_{R>0}\left(\int\limits_{Q_{n}}\left|\Delta_{R}(x)\right|^{p} dx\right)^{1/p}\leq A p,1\leq p<\infty,
\]

where

\[
\Delta_{R}(x)=K_{R}^{(n-1)/2}(x)-\pi^{(n-1)/2}\Gamma\biggl(\frac{n+1}{2}\biggr)R^{1/2}\big|x\big|^{-n+1/2}J_{n-1/2}(2\pi R|x|\big).
\]

For details see Stein [7].

\frp{6}{10} The following convexity theorem for Riesz means of numerical series can be used in the proof of Theorem \ref{thm:c7-4-3}. Consider a numerical series $\sum_{k \geq 0} c_k$ and its Riesz means $\sigma_R^\alpha = \sum_{0 \leq k < R} (1 - k/R)^\alpha c_k, \alpha \geq 0$. Suppose that $\sigma_R^\alpha = O(R^{\alpha_t})$ as $R \to \infty$, for $j = 0, 1$. Then, if $0 < \theta < 1$, $\alpha = \alpha_0(1 - \theta) + \alpha_1\theta$, and $a = a_0(1 - \theta) + a_1\theta$, we can conclude that $\sigma_R^\alpha = O(R^\alpha)$ as $R \to \infty$. See Riesz [1], and Chandrasekharan and Minukshandaran [1]. To apply this result in the proof of Theorem \ref{thm:c7-4-3} write

\[
\sigma_{R}^{\alpha}=K_{R}^{1/2}(x^{0})=\sum_{|{\bf m}|^{2}<R}\ \left(1-\frac{|{\bf m}|^{2}}{R}\right)^{\alpha}e^{2\pi i{\bf m}\cdot{\bf x}^{0}}
\]

for

\[
x_{0}\notin\Lambda\text{(where}c_{k}=\sum_{|m|^{2}=k}e^{2\pi i m\cdot x^{0}).}
\]

We know (see the inequality following \eqref{eq:c7-4-28}) that $\sigma_R^\alpha = O(R^{(1/2)[(n-1)/2-\alpha]})$ if $\alpha > (n-1)/2$. If it were true that $\sigma_R^0 = O(R^{(1/2)[(n-1)/2]}$) then the convexity theorem just quoted would imply that $\sigma_R^{(n-1)/2} = O(1)$, which contradicts Lemma \ref{lem:c7-4-6}, for almost every $x^0$.

\bibnoteshead

For general references dealing with the elementary aspects of multiple Fourier series and the Poisson summation formula, see the classical book of

Bochner [3] and also his later tract [6]. The significance of the critical index (see Corollary \ref{cor:c7-2-15}) goes back to Bochner's paper [7]. The n-dimensional form of Theorem \ref{thm:c7-2-17} may be found in Wainger's memoir [1].

The main results of the third section, Theorems \ref{thm:c7-3-8} and Corollary \ref{cor:c7-3-16} are due to de Leeuw [1]. The general ideas of the argument leading to Theorem \ref{thm:c7-3-18} are classical, but the present sharp form seems to be new. See also Igari [1].

The failure of localization at the critical index for $L^1(T_n)$ was proved in Bochner's paper [7] that was alluded to above. The existence of an $L^1(T_n)$ function $(n > 1)$ whose Fourier series is almost nowhere summable by the Riesz means at the critical index is contained in Stein [5]. The present stronger form of Theorem \ref{thm:c7-4-2}, however, is new; the same is true of Theorem \ref{thm:c7-4-3}. All of these results have as one of their starting points the technique introduced by Bochner, which involves the set $S$ of points $x$ where $\{|x - m|\}$ are linearly independent over the rationals. For Theorem \ref{thm:c7-5-1} see Stein [4]; also [2], and [7] which are closely related. The reader may also consult the Survey paper of Shapiro [1], which covers several topics in multiple Fourier series.

